% Options for packages loaded elsewhere
\PassOptionsToPackage{unicode}{hyperref}
\PassOptionsToPackage{hyphens}{url}
\documentclass[
  ukrainian,
]{article}
\usepackage{xcolor}
\usepackage[margin=2.5cm]{geometry}
\usepackage{amsmath,amssymb}
\setcounter{secnumdepth}{-\maxdimen} % remove section numbering
\usepackage{iftex}
\ifPDFTeX
  \usepackage[T1]{fontenc}
  \usepackage[utf8]{inputenc}
  \usepackage{textcomp} % provide euro and other symbols
\else % if luatex or xetex
  \usepackage{unicode-math} % this also loads fontspec
  \defaultfontfeatures{Scale=MatchLowercase}
  \defaultfontfeatures[\rmfamily]{Ligatures=TeX,Scale=1}
\fi
\usepackage{lmodern}
\ifPDFTeX\else
\setmainfont{DejaVu Serif}
\setsansfont{DejaVu Sans}
\setmonofont{DejaVu Sans Mono}
\setmathfont{STIX Two Math}
\fi
\ifPDFTeX\else
  % xetex/luatex font selection
\fi
% Use upquote if available, for straight quotes in verbatim environments
\IfFileExists{upquote.sty}{\usepackage{upquote}}{}
\IfFileExists{microtype.sty}{% use microtype if available
  \usepackage[]{microtype}
  \UseMicrotypeSet[protrusion]{basicmath} % disable protrusion for tt fonts
}{}
\makeatletter
\@ifundefined{KOMAClassName}{% if non-KOMA class
  \IfFileExists{parskip.sty}{%
    \usepackage{parskip}
  }{% else
    \setlength{\parindent}{0pt}
    \setlength{\parskip}{6pt plus 2pt minus 1pt}}
}{% if KOMA class
  \KOMAoptions{parskip=half}}
\makeatother
\usepackage{longtable,booktabs,array}
\newcounter{none} % for unnumbered tables
\usepackage{calc} % for calculating minipage widths
% Correct order of tables after \paragraph or \subparagraph
\usepackage{etoolbox}
\makeatletter
\patchcmd\longtable{\par}{\if@noskipsec\mbox{}\fi\par}{}{}
\makeatother
% Allow footnotes in longtable head/foot
\IfFileExists{footnotehyper.sty}{\usepackage{footnotehyper}}{\usepackage{footnote}}
\makesavenoteenv{longtable}
\usepackage{graphicx}
\graphicspath{{/Users/kalinicenkovitalijmikolajovic/.codex/worktrees/5c15/Script_for_translate_book/Output_Final/}}
\setkeys{Gin}{width=0.82\linewidth,height=0.82\textheight,keepaspectratio}
\makeatletter
\newsavebox\pandoc@box
\newcommand*\pandocbounded[1]{% scales image to fit in text height/width
  \sbox\pandoc@box{#1}%
  \Gscale@div\@tempa{\textheight}{\dimexpr\ht\pandoc@box+\dp\pandoc@box\relax}%
  \Gscale@div\@tempb{\linewidth}{\wd\pandoc@box}%
  \ifdim\@tempb\p@<\@tempa\p@\let\@tempa\@tempb\fi% select the smaller of both
  \ifdim\@tempa\p@<\p@\scalebox{\@tempa}{\usebox\pandoc@box}%
  \else\usebox{\pandoc@box}%
  \fi%
}
% Set default figure placement to htbp
\def\fps@figure{htbp}
\makeatother
\ifLuaTeX
\usepackage[bidi=basic,shorthands=off,provide=*]{babel}
\else
\usepackage[bidi=default,shorthands=off,provide=*]{babel}
\fi
\ifLuaTeX
  \usepackage{selnolig} % disable illegal ligatures
\fi
\setlength{\emergencystretch}{3em} % prevent overfull lines
\providecommand{\tightlist}{%
  \setlength{\itemsep}{0pt}\setlength{\parskip}{0pt}}
\usepackage{bookmark}
\IfFileExists{xurl.sty}{\usepackage{xurl}}{} % add URL line breaks if available
\urlstyle{same}
\hypersetup{
  pdflang={uk},
  hidelinks,
  pdfcreator={LaTeX via pandoc}}

\author{}
\date{}

\begin{document}

\pandocbounded{\includegraphics[keepaspectratio]{images/_page_150_Picture_5.jpeg}}

Рис. 20.3. Стоячі хвилі \(u_n(x, t) = X_n(x)T_n(t)\).

В нашей задаче уравнение и граничные условия линейны и однородны,
поэтому произвольная сумма элементарных стоячих волн также является
решением волнового уравнения и удовлетворяет граничным условиям. Мы
должны теперь найти такую суперпозицию стоячих волн, чтобы она
удовлетворяла также и начальным условиям, тогда получим решение нашей
задачи. Подстановка суммы

Зверніть увагу на зображення.

\pandocbounded{\includegraphics[keepaspectratio]{images/_page_154_Figure_5.jpeg}}

Головна різниця між поперечними коливаннями тонкої балки та поперечними
коливаннями струни полягає в тому, що промінь чинить опір згину. Не
заглиблюючись у механіку тонких балок, припустимо, що, враховуючи опір
згинання (замість рівняння хвилі) веде до рівняння четвертого порядку

(21.1)

\section{Свободно опирающаяся
балка}\label{ux441ux432ux43eux431ux43eux434ux43dux43e-ux43eux43fux438ux440ux430ux44eux449ux430ux44fux441ux44f-ux431ux430ux43bux43aux430}

Розглянемо невеликі коливання тонкої балки, кінці якої вільно лежать на
двох опорах. Коли ми кажемо «вільно відпочивати», ми маємо на увазі, що
кінці балки не рухаються, але

Клони пучка на кінцевих точках можна змінювати (кінці пучка фіксуються
за допомогою штифтових пристроїв, рис. 21.1).

Совершенно ясно, что на концах балки должны выполняться граничные
условия

\[
F(\omega) = \frac{?}{\pi} \int_{0}^{\infty} f(x) \sin(\omega x) dx
\]

\(u(1, t) = 0,\) але зовсім не очевидно, що на кінцях балки мають бути
виконані ще дві граничні умови:

\[
0 < \omega < x
\]

\(u_{xx}(1, t) = 0.\) Використовуючи теорію тонких балок, можна
показати, що момент згинання в балці пропорційний значенню \(u_{xx}\) ,
а момент згинання на вільно лежачому кінці має дорівнювати P I C

\pandocbounded{\includegraphics[keepaspectratio]{images/_page_155_Figure_8.jpeg}}

.21.1. Свободно оппрающаяся балка.

Нуль вен. Відповідно, зображення, зображене на рис. 21.1 Осцилюючий
пучок описується такою змішаною задачею (для простоти значення o
вважається рівним одиниці):

(21.2)

f'\,'(0)

\begin{enumerate}
\def\labelenumi{\arabic{enumi}.}
\setcounter{enumi}{1}
\tightlist
\item
  \(f(\alpha x)\)
\end{enumerate}

\[
\begin{aligned} & (\text{Y} \cdot \Pi) \quad u_{tt} = -u_{xxxx}, \quad 0 < x < 1, \quad 0 < t < \infty, \\ & \begin{cases} u \cdot (0, \ t) = 0, \\ u_{xx} \cdot (0, \ t) = 0, \\ u \cdot (1, \ t) = 0, \\ u_{xx} \cdot (1, \ t) = 0, \end{cases} \quad 0 < t < \infty, \\ & \begin{cases} u \cdot (x, \ 0) = f(x), \\ u_{1}(x, \ 0) = g(x). \end{cases} \quad 0 \leqslant x \leqslant 1, \end{aligned}
\]

\begin{enumerate}
\def\labelenumi{\arabic{enumi}.}
\setcounter{enumi}{2}
\item
  \(e^{-\alpha x}\)
\item
\end{enumerate}

Щоб розв'язати цю задачу, ми використаємо метод розділення змінних. Ми
будемо шукати лише періодичні розв'язки, тобто флуктуації

(21.3)

\[
[2/\pi\omega]^{1/2}
\]

Зверніть увагу, що вибір розв'язку вигляду (21.3) фактично означає, що
ми взяли константу розділення у методі розділення змінних як від'ємну.

Подставим теперь (21.3) в уравнение колебаний балки и получим уравнение
для \(X\left(x\right)\)

\[
\frac{2}{\pi\omega} [1-\cos(\omega a)]
\]

общее решение которого записывается в виде

1

Щоб визначити коефіцієнти C, D, E та F, підставляємо загальний розв'язок
у граничні умови

\[
e^{-1}a + \omega
\frac{1}{2}e^{-\alpha}\sin\omega
\begin{split} u\left(0,\ t\right) &= 0 \Rightarrow X\left(0\right)T\left(t\right) = 0 \Rightarrow X\left(0\right) = 0 \Rightarrow C + E = 0 \\ u_{xx}\left(0,\ t\right) &= 0 \Rightarrow X''\left(0\right)T\left(t\right) = 0 \Rightarrow X''\left(0\right) = 0 \Rightarrow -C + E = 0 \\ u\left(1,\ t\right) &= 0 \Rightarrow D\sin\sqrt{\omega} + F\sin\sqrt{\omega} = 0, \\ u_{xx}\left(1,\ t\right) &= 0 \Rightarrow -D\sin\sqrt{\omega} + F\sin\sqrt{\omega} = 0. \end{split}
\frac{1-e^{-a\omega}}{\omega}
\frac{2}{\pi}F''(\omega)
F \sinh \sqrt{\omega} = 0,
f(x) = \int_{0}^{\infty} F(\omega) \cos(\omega x) d\omega \qquad \qquad F(x) = \frac{2}{\pi} \int_{0}^{\infty} f(x) \cos(\omega x) dx
D \sin \sqrt{\omega} = 0.
\]

\begin{enumerate}
\def\labelenumi{\arabic{enumi}.}
\tightlist
\item
\end{enumerate}

\[
F = 0
\]

\section{ПЕРЕХОД К БЕЗРАЗМЕРНЫМ ПЕРЕМЕННЫМ МЕТА ЛЕКЦІЇ: Показать, каким
образом краевые, начальные и другие виды физических задач можно записать
в безразмерном виде. В конкретных областях знания (физике, химии,
биологии или экономике) уравнения записываются
по-разному.}\label{ux43fux435ux440ux435ux445ux43eux434-ux43a-ux431ux435ux437ux440ux430ux437ux43cux435ux440ux43dux44bux43c-ux43fux435ux440ux435ux43cux435ux43dux43dux44bux43c-ux43cux435ux442ux430-ux43bux435ux43aux446ux456ux457-ux43fux43eux43aux430ux437ux430ux442ux44c-ux43aux430ux43aux438ux43c-ux43eux431ux440ux430ux437ux43eux43c-ux43aux440ux430ux435ux432ux44bux435-ux43dux430ux447ux430ux43bux44cux43dux44bux435-ux438-ux434ux440ux443ux433ux438ux435-ux432ux438ux434ux44b-ux444ux438ux437ux438ux447ux435ux441ux43aux438ux445-ux437ux430ux434ux430ux447-ux43cux43eux436ux43dux43e-ux437ux430ux43fux438ux441ux430ux442ux44c-ux432-ux431ux435ux437ux440ux430ux437ux43cux435ux440ux43dux43eux43c-ux432ux438ux434ux435.-ux432-ux43aux43eux43dux43aux440ux435ux442ux43dux44bux445-ux43eux431ux43bux430ux441ux442ux44fux445-ux437ux43dux430ux43dux438ux44f-ux444ux438ux437ux438ux43aux435-ux445ux438ux43cux438ux438-ux431ux438ux43eux43bux43eux433ux438ux438-ux438ux43bux438-ux44dux43aux43eux43dux43eux43cux438ux43aux435-ux443ux440ux430ux432ux43dux435ux43dux438ux44f-ux437ux430ux43fux438ux441ux44bux432ux430ux44eux442ux441ux44f-ux43fux43e-ux440ux430ux437ux43dux43eux43cux443.}

Переход к безразмерным переменным позволяет привести уравнения к одному
и тому же виду. Именно по этой причине при изучении уравнений с частными
производными в математике стараются отвлекаться от физического смысла
входящих в них параметров. Вот почему химик, физик или биолог должен
преобразовать свое уравнение к той форме, которая принята в нашей книге.

Основная идея введения новых (безразмерных) переменных состоит в том,
что после перехода к безразмерным переменным вадача становится чисто
математической и уже не содержит

карактерных физических констант. Именно таким способом многие различные
уравнения физики, биологии и химии, содержащие всякие нюансы, связанные
с физическими параметрами, приводятся к одной и той же простой форме
(см. рис. 22.1).

Познакомимся с этим преобразованием на простом примере.

\pandocbounded{\includegraphics[keepaspectratio]{images/_page_161_Figure_4.jpeg}}

Рис. 22.1. Некоторые задачи, приводящиеся к одной и той же безразмерной
форме.

\section{Приведение диффузионной задачи к безразмерному
виду}\label{ux43fux440ux438ux432ux435ux434ux435ux43dux438ux435-ux434ux438ux444ux444ux443ux437ux438ux43eux43dux43dux43eux439-ux437ux430ux434ux430ux447ux438-ux43a-ux431ux435ux437ux440ux430ux437ux43cux435ux440ux43dux43eux43cux443-ux432ux438ux434ux443}

\begin{verbatim}
Начнем со смешанной задачи для отрезка с начальным условием на температуру  u\left(x,\,0\right)=\sin\left(\pi x/L\right)  и граничными температурами на концах  T_1  и  T_2 . Другими словами, мы решаем задачу
\end{verbatim}

(22.1)

\begin{verbatim}
\begin{array}{ccccccccccccccccccccccccccccccccccc
\end{verbatim}

(см. рис. 22.2).

\pandocbounded{\includegraphics[keepaspectratio]{images/_page_161_Figure_10.jpeg}}

Рис. 22.2. Задача теплопроводности в плоскости переменных (х, т).

Наша цель --- дать новую эквивалентную (22.1) постановку задачи таким
образом, чтобы

в новой задаче не было никаких физических параметров (вроде α);

\begin{enumerate}
\def\labelenumi{\arabic{enumi})}
\setcounter{enumi}{1}
\tightlist
\item
  начальные и граничные условия стали проще.
\end{enumerate}

\begin{verbatim}
Для осуществления этой цели, мы введем три новых безразмерных переменных U,   \xi   и   \tau  , чтобы заменить прежние u, x и t по схеме
\end{verbatim}

\[
u \to U
\]

\begin{verbatim}
(безразмерная температура),   x \to \xi   (безразмерная длина),   t \to \tau   (безразмерное время).
\end{verbatim}

Мы проведем все три преобразования одновременно.

\begin{verbatim}
# Преобразование зависимой переменной  u \rightarrow U
\end{verbatim}

Определим функцию \(U(x, t)\) по формуле

\[
U(x, t) = \frac{u(x, t) - T_1}{T_2 - T_1}
\]

.

\begin{verbatim}
Ясно, что новая температура   U\left(x,\,t\right)   безразмерна, поскольку мы делим градусы Цельсия на градусы Цельсия. Новые граничные условия для   U\left(x,\,t\right)   при x=0 и x=L примут вид   U\left(0,\,t\right)=0   и   U\left(L,\,t\right)=1  . Давайте поставим задачу для новой функции   U\left(x,\,t\right)  . После несложных преобразований задача (22.1) переходит в новую задачу
\end{verbatim}

задачу

\begin{verbatim}
\begin{array}{ll} \text{(УЧП)} & U_t = \alpha^2 U_{xx}, \quad 0 < x < L, \quad 0 < t < \infty, \\ U_t = 0, \quad 0 < t < \infty, \\ U_t = 0, \quad 0 < t < \infty, \\ U_t = 0, \quad 0 < t < \infty, \\ U_t = 0, \quad 0 < t < \infty, \\ U_t = 0, \quad 0 < t < \infty, \\ U_t = 0, \quad 0 < t < \infty, \\ U_t = 0, \quad 0 < t < \infty, \\ U_t = 0, \quad 0 < t < \infty, \\ U_t = 0, \quad 0 < t < \infty, \\ U_t = 0, \quad 0 < t < \infty, \\ U_t = 0, \quad 0 < t < \infty, \\ U_t = 0, \quad 0 < t < \infty, \\ U_t = 0, \quad 0 < t < \infty, \\ U_t = 0, \quad 0 < t < \infty, \\ U_t = 0, \quad 0 < t < \infty, \\ U_t = 0, \quad 0 < t < \infty, \\ U_t = 0, \quad 0 < t < \infty, \\ U_t = 0, \quad 0 < t < \infty, \\ U_t = 0, \quad 0 < t < \infty, \\ U_t = 0, \quad 0 < t < \infty, \\ U_t = 0, \quad 0 < t < \infty, \\ U_t = 0, \quad 0 < t < \infty, \\ U_t = 0, \quad 0 < t < \infty, \\ U_t = 0, \quad 0 < t < \infty, \\ U_t = 0, \quad 0 < t < \infty, \\ U_t = 0, \quad 0 < t < \infty, \\ U_t = 0, \quad 0 < t < \infty, \\ U_t = 0, \quad 0 < t < \infty, \\ U_t = 0, \quad 0 < t < \infty, \\ U_t = 0, \quad 0 < t < \infty, \\ U_t = 0, \quad 0 < t < \infty, \\ U_t = 0, \quad 0 < t < \infty, \\ U_t = 0, \quad 0 < t < \infty, \\ U_t = 0, \quad 0 < t < \infty, \\ U_t = 0, \quad 0 < t < \infty, \\ U_t = 0, \quad 0 < t < \infty, \\ U_t = 0, \quad 0 < t < \infty, \\ U_t = 0, \quad 0 < t < \infty, \\ U_t = 0, \quad 0 < t < \infty, \\ U_t = 0, \quad 0 < t < \infty, \\ U_t = 0, \quad 0 < t < \infty, \\ U_t = 0, \quad 0 < t < \infty, \\ U_t = 0, \quad 0 < t < \infty, \\ U_t = 0, \quad 0 < t < \infty, \\ U_t = 0, \quad 0 < t < \infty, \\ U_t = 0, \quad 0 < t < \infty, \\ U_t = 0, \quad 0 < t < \infty, \\ U_t = 0, \quad 0 < t < \infty, \\ U_t = 0, \quad 0 < t < \infty, \\ U_t = 0, \quad 0 < t < \infty, \\ U_t = 0, \quad 0 < t < \infty, \\ U_t = 0, \quad 0 < t < \infty, \\ U_t = 0, \quad 0 < t < \infty, \\ U_t = 0, \quad 0 < t < \infty, \\ U_t = 0, \quad 0 < t < \infty, \\ U_t = 0, \quad 0 < t < \infty, \\ U_t = 0, \quad 0 < t < \infty, \\ U_t = 0, \quad 0 < t < \infty, \\ U_t = 0, \quad 0 < t < \infty, \\ U_t = 0, \quad 0 < t < \infty, \\ U_t = 0, \quad 0 < t < \infty, \\ U_t = 0, \quad 0 < t < \infty, \\ U_t = 0, \quad 0 < t < \infty, \\ U_t = 0, \quad 0 < t < \infty, \\ U_t = 0, \quad 0 < t < \infty, \\ U_t = 0, \quad 0 < t < \infty, \\ U_t = 0, \quad 0 < t < \infty, \\ U_t = 0, \quad 0 < t < \infty, \\ U_t = 0, \quad 0 < t < \infty, \\ U_t = 0, \quad 0 < t < \infty, \\ U_t = 0, \quad 0 < t < \infty, \\ U_t = 0, \quad 0 < t < \infty, \\ U_t = 0, \quad 0 < t < \infty, \\ U_t = 0, \quad 0 < t < \infty, \\ U_t = 0, \quad 0 < t < \infty, \\ U_t = 0, \quad 0 < t < \infty, \\ U_t = 0, \quad 0 < t < \infty, \\ U_t = 0, \quad 0 < t < \infty, \\ U_t = 0, \quad 0 < t < \infty, \\ U_t = 0, \quad 0 < t < \infty, \\ U_t = 0, \quad 0 < t < \infty, \\ U_t = 0, \quad 0 < t < \infty, \\ U_t
\end{verbatim}

\begin{verbatim}
Если угодно, то здесь можно остановиться, решить задачу для   U\left(x,\, t\right)   и затем найти   u\left(x,\,t\right)   по формуле
\end{verbatim}

\[
u(x, t) = T_1 + (T_2 - T_1) U(x, t).
\]

Однако мы пойдем дальше и займемся преобразованием независимых
переменных x и t.

\begin{verbatim}
# Преобразование пространственной координаты  x \rightarrow \xi
\end{verbatim}

\begin{verbatim}
Представляется совершенно очевидным, как следует выбрать безразмерную переменную   \xi  . Поскольку   0 \leqslant x \leqslant L  , полагаем
\end{verbatim}

\[
\xi = x/L
\]

.

После вычисления производных

\[
U_x = U_{\xi} \xi_x = \frac{1}{L} U_{\xi},
U_{xx} = \frac{1}{L^2} U_{\xi\xi}
\]

наступна задача стає очевидною (у змінних U, \(\xi\) , t):

\[
\begin{array}{ll} (\text{YM}\Pi) & U_t = (\alpha/L)^2 \, U_{\xi\xi}, \quad 0 < \xi < 1, \quad 0 < t < \infty, \\ (22.3) & (\Gamma \text{Y}) & \begin{cases} U(0,\ t) = 0, \\ U(1,\ t) = 1, \\ U(\xi,\ 0) = \frac{\sin\left(\pi\xi\right) - T_1}{T_2 - T_1}, \quad 0 \leqslant \xi \leqslant 1. \end{cases} \end{array}
\]

Ми вже пройшли дві третини шляху. Останній крок пов'язаний із введенням
безрозмірного часу \(\tau\) для виключення коефіцієнта \([\alpha/L]^2\)
з диференціального рівняння.

\section{\texorpdfstring{\(t \rightarrow \tau\) Перетворення
часу}{t \textbackslash rightarrow \textbackslash tau Перетворення часу}}\label{t-rightarrow-tau-ux43fux435ux440ux435ux442ux432ux43eux440ux435ux43dux43dux44f-ux447ux430ux441ux443}

Як обрати час без розмірів? Це не так очевидно, як у випадку перших двох
змінних. Однак, оскільки наша мета --- виключити сталу \([\alpha/L]^2\)
з рівняння (22.3), ми діятимемо наступним чином.

\begin{itemize}
\item
  \begin{enumerate}
  \def\labelenumi{\arabic{enumi}.}
  \tightlist
  \item
    Спробуємо замінити вигляд \(\tau = ct\) , де c --- невідома
    константа.
  \end{enumerate}

  \begin{itemize}
  \tightlist
  \item
    \begin{enumerate}
    \def\labelenumi{\arabic{enumi}.}
    \setcounter{enumi}{1}
    \tightlist
    \item
      Давайте розрахуємо \(u_t = u_{\tau} \tau_t = c u_{\tau}\) .
    \end{enumerate}
  \item
    \begin{enumerate}
    \def\labelenumi{\arabic{enumi}.}
    \setcounter{enumi}{2}
    \tightlist
    \item
      Підставимо цю похідну у рівняння (22.3) і отримаємо
    \end{enumerate}
  \end{itemize}
\end{itemize}

\[
cu_{\mathbf{t}} = [\alpha/L]^2 u_{\xi\xi}.
\]

Тому потрібно обрати \(c = [\alpha/L]^2\) , тобто новий час слід вводити
за формулою

\[
\tau = [\alpha/L]^2 t.
\]

Застосовуючи це перетворення до попередньої задачі (22.3), ми можемо
повністю перейти до безрозмірних змінних U, \(\xi\) та \(\tau_i\)

\begin{verbatim}
\begin{array}{cccccccccccccccccccccccccccccccccccc
\end{verbatim}

Новая безразмерная система обладает следующими свойствами:

\begin{enumerate}
\def\labelenumi{\arabic{enumi}.}
\item
  Уравнение не содержит параметров.
\item
  Граничные условия имеют простой вид.
\item
  Початкова умова суттєво не змінилася (тобто містить лише відомі
  функції).
\item
  Завдання стало простішим і компактнішим порівняно з оригінальним.
  Розв'язок цієї проблеми можна побудувати раз і назавжди, тож якщо
  читач перетворює свою початкову задачу (22.1) у безрозмірну форму
  (22.4) і знаходить її розв'язок \(U(\xi, \tau)\) у підручнику чи
  науковому журналі, то він може знайти розв'язок початкової задачі
  (22.1) простим перерахуванням
\end{enumerate}

\[
u(x, t) = T_1 + (T_2 - T_1) U(x/L, \alpha^2 t/L^2).
\]

На цьому завершується наша дискусія проблеми перетворення задач у
безрозмірний вигляд. Не існує універсальних правил, за якими слід
вводити нові змінні, тут слід покладатися на фізичну інтуїцію і
пробувати різні варіанти.

Завершимо цю лекцію простим прикладом перетворення задачі на
безрозмірний вигляд, розв'язання нової задачі і повернення до початкових
змінних.

\section{Приклад перетворення гіперболічної задачі на безрозмірний
вигляд}\label{ux43fux440ux438ux43aux43bux430ux434-ux43fux435ux440ux435ux442ux432ux43eux440ux435ux43dux43dux44f-ux433ux456ux43fux435ux440ux431ux43eux43bux456ux447ux43dux43eux457-ux437ux430ux434ux430ux447ux456-ux43dux430-ux431ux435ux437ux440ux43eux437ux43cux456ux440ux43dux438ux439-ux432ux438ux433ux43bux44fux434}

Рассмотрим колеблющуюся струну

(22.5)

\[
\begin{array}{llll} (\text{УЧП}) & u_{tt} = \alpha^2 u_{xx}, & 0 < x < L, & 0 < t < \infty, \\ u(0, t) = 0, & 0 < t < \infty, \\ u(L, t) = 0, & 0 < t < \infty, \\ u(X, 0) = \sin(\pi x/L) + (0.5)\sin(3\pi x/L), \\ u_t(x, 0) = 0. & \end{array}
\]

Перетворюючи пояснювальні змінні у безрозмірний вигляд (немає потреби
трансформувати вас)

\[
\xi = x/L
\]

, \(\tau = [\alpha/L] t\) получаем новую задачу

\[
(\text{YYII}) \qquad u_{\tau\tau} = u_{\xi\xi}, \quad 0 < \xi < 1, \quad 0 < \tau < \infty, 1
(\text{\GammaY}) \quad
\begin{cases} u(0, \tau) = 0, \\ u(1, \tau) = 0, \\ u(\xi, 0) = \sin(\pi\xi) + 0.5\sin(3\pi\xi), \\ u_{\xi}(\xi, 0) = 0, \end{cases}
\quad 0 < \xi \leqslant 1,
(\text{HY}) \quad
\begin{cases} u(\xi, 0) = \sin(\pi\xi) + 0.5\sin(3\pi\xi), \\ u_{\xi}(\xi, 0) = 0, \end{cases}
\quad 0 \leqslant \xi \leqslant 1,
\]

решением которой является функция

\[
u(\xi, \tau) = \cos(\pi \tau) \sin(\pi \xi) + 0.5 \cos(3\pi \tau) \sin(3\pi \xi).
\]

Якщо повернутися до координат x і t, розв'язок задачі (22.5) набуває
вигляду

\[
u(x, t) = \cos(\pi \alpha t/L)\sin(\pi x/L) + 0.5\cos(3\pi \alpha t/L)\sin(3\pi x/L).
\]

\section{ПОВАГА - 1. Аналіз розмірності особливо важливий у чисельних
методах, оскільки більшість програм написані для розв'язання загальних
математичних задач і не передбачають використання великої кількості
фізичних параметрів. Ті, хто хоче користуватися цими програмами, повинні
перетворити свою задачу у відповідну форму, розв'язати її на комп'ютері,
а потім довести числові результати до потрібного
розміру.}\label{ux43fux43eux432ux430ux433ux430---1.-ux430ux43dux430ux43bux456ux437-ux440ux43eux437ux43cux456ux440ux43dux43eux441ux442ux456-ux43eux441ux43eux431ux43bux438ux432ux43e-ux432ux430ux436ux43bux438ux432ux438ux439-ux443-ux447ux438ux441ux435ux43bux44cux43dux438ux445-ux43cux435ux442ux43eux434ux430ux445-ux43eux441ux43aux456ux43bux44cux43aux438-ux431ux456ux43bux44cux448ux456ux441ux442ux44c-ux43fux440ux43eux433ux440ux430ux43c-ux43dux430ux43fux438ux441ux430ux43dux456-ux434ux43bux44f-ux440ux43eux437ux432ux44fux437ux430ux43dux43dux44f-ux437ux430ux433ux430ux43bux44cux43dux438ux445-ux43cux430ux442ux435ux43cux430ux442ux438ux447ux43dux438ux445-ux437ux430ux434ux430ux447-ux456-ux43dux435-ux43fux435ux440ux435ux434ux431ux430ux447ux430ux44eux442ux44c-ux432ux438ux43aux43eux440ux438ux441ux442ux430ux43dux43dux44f-ux432ux435ux43bux438ux43aux43eux457-ux43aux456ux43bux44cux43aux43eux441ux442ux456-ux444ux456ux437ux438ux447ux43dux438ux445-ux43fux430ux440ux430ux43cux435ux442ux440ux456ux432.-ux442ux456-ux445ux442ux43e-ux445ux43eux447ux435-ux43aux43eux440ux438ux441ux442ux443ux432ux430ux442ux438ux441ux44f-ux446ux438ux43cux438-ux43fux440ux43eux433ux440ux430ux43cux430ux43cux438-ux43fux43eux432ux438ux43dux43dux456-ux43fux435ux440ux435ux442ux432ux43eux440ux438ux442ux438-ux441ux432ux43eux44e-ux437ux430ux434ux430ux447ux443-ux443-ux432ux456ux434ux43fux43eux432ux456ux434ux43dux443-ux444ux43eux440ux43cux443-ux440ux43eux437ux432ux44fux437ux430ux442ux438-ux457ux457-ux43dux430-ux43aux43eux43cux43fux44eux442ux435ux440ux456-ux430-ux43fux43eux442ux456ux43c-ux434ux43eux432ux435ux441ux442ux438-ux447ux438ux441ux43bux43eux432ux456-ux440ux435ux437ux443ux43bux44cux442ux430ux442ux438-ux434ux43e-ux43fux43eux442ux440ux456ux431ux43dux43eux433ux43e-ux440ux43eux437ux43cux456ux440ux443.}

\begin{itemize}
\tightlist
\item
  \begin{enumerate}
  \def\labelenumi{\arabic{enumi}.}
  \setcounter{enumi}{1}
  \tightlist
  \item
    Розмірний аналіз дозволяє математикам працювати з диференціальним
    рівнянням у похідних похідних, не виключаючи багатьох параметрів і
    констант, які не впливають на математичну суть задачі.
  \end{enumerate}
\item
  \begin{enumerate}
  \def\labelenumi{\arabic{enumi}.}
  \setcounter{enumi}{2}
  \tightlist
  \item
    Не завжди потрібно конвертувати всі змінні у безрозмірний вигляд,
    іноді достатньо перетворити одну або дві змінні. \# ЗАВДАННЯ
  \end{enumerate}
\end{itemize}

\begin{enumerate}
\def\labelenumi{\arabic{enumi}.}
\tightlist
\item
  Шляхом заміни змінних
\end{enumerate}

\[
\xi = x/L
\]

, \(\tau = [\alpha/L] t\) преобразуйте задачу (22.5) к безразмерному
виду (22.6).

\begin{enumerate}
\def\labelenumi{\arabic{enumi}.}
\setcounter{enumi}{1}
\tightlist
\item
  Найдите безразмерную форму для задачи
\end{enumerate}

\begin{verbatim}
\begin{array}{lll} (\text{Y} \, \Pi) & u_t = \alpha^2 u_{xx}, & 0 < x < L, \\ (\Gamma \, \text{Y}) & \left\{ \begin{array}{ll} u \, (0, \, t) = T_1, \\ u \, (L, \, t) = 0, \\ u \, (x, \, 0) = T_2, & 0 \leqslant x \leqslant L. \end{array} \end{array}
\end{verbatim}

\begin{verbatim}
\right. \\ \end{array}
\end{verbatim}

\begin{enumerate}
\def\labelenumi{\arabic{enumi}.}
\setcounter{enumi}{2}
\tightlist
\item
  С помощью замены
\end{enumerate}

\[
U(x, t) = \frac{u(x, t) - T_1}{T_2 - T_1}
\]

преобразуйте задачу (22.1) в задачу (22.2).

\begin{verbatim}
4. Почему замена переменной   \tau = \alpha t   приводит к исключению параметра   \alpha   из волнового уравнения   u_{tt} = \alpha^2 u_{xx}  ? В чем физический смысл этой замены? Напоминаем вам, что   \alpha  —скорость
\end{verbatim}

распространения волн. При поиске новых координат в большинстве случаев
главную роль играет интуиция.

Как выбрать новую пространственную координату ξ, чтобы исключить и из
уравнения

\[
u_t + vu_x = 0?
\]

Лекция 23

\section{КЛАССИФИКАЦИЯ УРАВНЕНИЙ С ЧАСТНЫМИ ПРОИЗВОДНЫМИ (КАНОНИЧЕСКАЯ
ФОРМА ГИПЕРБОЛИЧЕСКОГО
УРАВНЕНИЯ)}\label{ux43aux43bux430ux441ux441ux438ux444ux438ux43aux430ux446ux438ux44f-ux443ux440ux430ux432ux43dux435ux43dux438ux439-ux441-ux447ux430ux441ux442ux43dux44bux43cux438-ux43fux440ux43eux438ux437ux432ux43eux434ux43dux44bux43cux438-ux43aux430ux43dux43eux43dux438ux447ux435ux441ux43aux430ux44f-ux444ux43eux440ux43cux430-ux433ux438ux43fux435ux440ux431ux43eux43bux438ux447ux435ux441ux43aux43eux433ux43e-ux443ux440ux430ux432ux43dux435ux43dux438ux44f}

МЕТА ЛЕКЦІЇ: Показать, что линейное уравнение второго порядка с двумя
независимыми переменными

\[
Au_{xx} + Bu_{xy} + Cu_{yy} + Du_x + Eu_y + Fu = G
\]

(A, B, C, D, E, F и G являются функциями x и y или могут быть
константами) относится к одному из следующих типов:

\begin{itemize}
\tightlist
\item
  \begin{enumerate}
  \def\labelenumi{\arabic{enumi})}
  \tightlist
  \item
    гиперболическому (если B\^{}2 4AC \textgreater{} 0 ),
  \end{enumerate}
\item
  \begin{enumerate}
  \def\labelenumi{\arabic{enumi})}
  \setcounter{enumi}{1}
  \tightlist
  \item
    параболическому (если B\^{}2 4AC = ),
  \end{enumerate}
\item
  \begin{enumerate}
  \def\labelenumi{\arabic{enumi})}
  \setcounter{enumi}{2}
  \tightlist
  \item
    эллиптическому (если 0 B\^{}2 4AC \textless{} 0 ).
  \end{enumerate}
\end{itemize}

\begin{verbatim}
Показать, какие новые переменные   \xi = \xi(x, y)   и   \eta = \eta(x, y)   можно ввести вместо x и y, чтобы упростить исходное уравнение. В новых переменных   \xi   и   \eta   уравнение с частными производными приводится к одному из следующих видов (в вависимости от того, будет ли величина   B^2 - 4AC   положительна, равна нулю или отрицательна):
\end{verbatim}

\begin{itemize}
\tightlist
\item
  \begin{enumerate}
  \def\labelenumi{\arabic{enumi}.}
  \tightlist
  \item
  \end{enumerate}
\end{itemize}

\[
\begin{cases} u_{\xi\xi} u_{\eta\eta} = \Psi(\xi,  \eta, u, u_{\xi}, u_{\eta}) і \text{(дві канонічні форми гіперболічного рівняння),} \end{cases}
\]

\begin{verbatim}
- 2.   u_{\eta\eta} = \Phi(\xi, \eta, u, u_{\xi}, u_{\eta})   (каноническая форма параболического уравнения),
\end{verbatim}

\begin{verbatim}
- 3.   u_{\xi\xi} + u_{\eta\eta} = \Phi(\xi, \eta, u, u_{\xi}, u_{\eta})   (каноническая форма эллиптического уравнения),
\end{verbatim}

\begin{verbatim}
где   \Phi   и   \Psi  —функции от частных производных первого порядка, зависимой переменной u и новых независимых переменных   \xi   и   \eta  . Конкретный вид функций   \Phi   и   \Psi   зависит от исходного уравнения.
\end{verbatim}

У читателя может возникнуть мысль, что глава, посвященная классификации
уравнений, должна размещаться в самом начале книги. Вероятно, это
правильно, и многие книги действительно начинаются с этого материала.
Однако правда и то, что многие студенты не могут с энтузиазмом изучать
нечто такое, о чем у них нет никакого представления. Именно по этой
причине мы отложили изучение темы, посвященной классификации уравнений с
частными производными, до настоящего момента.

Наша задача - классифицировать все уравнения вида

(23 1)

\[
Au_{xx} + Bu_{xy} + Cu_{yy} + Du_x + Eu_y + Fu = G,
\]

где A, B, C, D, E, F и G---вообще говоря, функции {[}x и y. Согласно
предлагаемой схеме классификации уравнение называется

\begin{enumerate}
\def\labelenumi{\arabic{enumi})}
\item
  гиперболическим в точке (\(x_0\), \(y_0\)) , если B\^{}2(\(x_0\),
  \(y_0\)) - 4A(\(x_0\), \(y_0\)) \помножити на \помножити на C(\(x_0\),
  \(y_0\)) \textgreater{} 0 2) параболнческим в точке (\(x_0\), \(y_0\))
  , если B\^{}2(\(x_0\), \(y_0\)) - 4A(\(x_0\), \(y_0\)) \помножити на
  C(\(x_0\), \(y_0\)) = 0 ,
\item
  эллиптическим в точке (\(x_0\), \(y_0\)) , если B\^{}2(\(x_0\),
  \(y_0\)) - 4A(\(x_0\), \(y_0\)) \помножити на \помножити на C(\(x_0\),
  \(y_0\)) \textless{} 0 и в зависимости от типа его можно привести к
  одной из канонических форм. Для того чтобы было легче разобраться в
  классификационной схеме, мы приведем четыре примера гиперболических,
  параболических и эллиптических уравнений.
\end{enumerate}

\section{Примеры гиперболических, параболических и эллиптических
уравнений
10}\label{ux43fux440ux438ux43cux435ux440ux44b-ux433ux438ux43fux435ux440ux431ux43eux43bux438ux447ux435ux441ux43aux438ux445-ux43fux430ux440ux430ux431ux43eux43bux438ux447ux435ux441ux43aux438ux445-ux438-ux44dux43bux43bux438ux43fux442ux438ux447ux435ux441ux43aux438ux445-ux443ux440ux430ux432ux43dux435ux43dux438ux439-10}

\begin{enumerate}
\def\labelenumi{\arabic{enumi}.}
\tightlist
\item
  У равнение теплопроводности \(u_t = u_{xx}\)--- линейное уравнение
  второго порядка вида (23.1) с коэффициентами
\end{enumerate}

\[
A=1
,   B=0  ,   C=0  ,   D=0  ,   E=-1  ,   F=0  ,   G=0  ,
\]

так что B\^{}2-4AC=0 при всех x и t. Следовательно, это уравнение будет
параболическим при всех x и t. В общем уравнении (23 1) переменная y
играет роль времени t. Результат не изменится, если в общем уравнении
поменять местами x и y.

\begin{enumerate}
\def\labelenumi{\arabic{enumi}.}
\setcounter{enumi}{1}
\tightlist
\item
  Волновое уравнение \(u_{1t} = u_{xx}\)также принадлежит классу
  уравнений вида (23.1) с коэффициентами
\end{enumerate}

\[
A=1
,   B=0  ,   C=-1  ,   D=E=F=G=0  .
\]

Так как B\^{}2 - 4AC = 4 для всех x и t, то это уравнение является
еиперболическим при всех x и t.

1) Следует указать, что приведенные в примерах уравнения являются
частными случаями более общих параболических, гиперболических в
эллиптических уравнений, а основные их свойства характерны и для общего
случая.

\begin{enumerate}
\def\labelenumi{\arabic{enumi}.}
\setcounter{enumi}{2}
\tightlist
\item
  У равнение Лапласа \(u_{xx}\) + \(u_{yy} = 0\)является эллиптическим
  при всех x и y, так как B\^{}2 - 4AC = - 4 \textless{} 0 .
\end{enumerate}

\begin{verbatim}
4. Линейное уравнение   xu_{xx} + u_{yy} = \sin x   с переменными коэффициентами также принадлежит классу уравнений вида (23.1), но для него   B^2-4AC=-4x   и, следовательно, уравнение будет эллиптическим при x > 0,
\end{verbatim}

лараболическим при x=0, гиперболическим при x \textless{} 0.

Этот пример иллюстрирует тот факт, что тип уравнения о переменными
коэффициентами может изменяться от точки к точке.

Читатель должен обратить внимание на следующее: тип уравнения (23.1)
определяется только коэффициентами при вторых првизводных и никак не
зависит ни от коэффициентов при первых производных и самой функции, ни
от свободного члена.

\begin{verbatim}
Обратимся теперь к главной теме нашей лекции — приведению гиперболических уравнений к каноническому виду. Известно, что если в данной области пространства уравнение является гиперболическим, то вместо x и y можно ввести координаты   \xi   и   \eta  (характеристические переменные) таким образом, что уравнение примет простейшую форму
\end{verbatim}

(23.2)

\[
u_{\xi\eta} = \Phi(\xi, \eta, u, u_{\xi}, u_{\eta}).
\]

\begin{verbatim}
Это уравнение содержит только одну производную второго порядка   u_{\xi\eta}  , а функция   \Phi(\xi,) \eta,u,u_{\xi},u_{\eta})   зависит только от новых независимых переменных § и η, зависимой переменной и и первых производных из и ип. Конкретный вид функции Ф зависит, конечно, от исходного уравнения и формул перехода к новым координатам § и п. Научимся отыскивать эти функции.
\end{verbatim}

\section{Каноническая форма гиперболического
уравнения}\label{ux43aux430ux43dux43eux43dux438ux447ux435ux441ux43aux430ux44f-ux444ux43eux440ux43cux430-ux433ux438ux43fux435ux440ux431ux43eux43bux438ux447ux435ux441ux43aux43eux433ux43e-ux443ux440ux430ux432ux43dux435ux43dux438ux44f}

Рассмотрим уравнение общего вида

(23.3)

\[
Au_{xx} + Bu_{xy} + Cu_{yy} + Du_x + Eu_y + Fu = G
\]

и пусть в интересующей нас области B\^{}2 - 4AC \textgreater{} 0 . Нам
нужно ввести новые координаты 1)

\[
\xi = \xi(x, y),
\eta = \eta(x, y)
\]

тому в рівнянні (23.3) залишається лише одна часткова похідна другого
порядку \(u_{\xi\eta}\) (не слід шукати заміну змінним, за допомогою
якої гіперболічне вирівнювання

1) Необхідно, щоб це перетворення координат було локально оборотним і
подвійно диференційованим. --- Ед. Ед.

У цьому випадку це неможливо).

Вычислим сначала частные производные

\[
\begin{array}{c} u_x = u_\xi \xi_x + u_\eta \eta_x, \\ u_y = u_\xi \xi_y + u_\eta \eta_y, \\ (23.4) \ u_{xx} = u_{\xi\xi} \xi_x^2 + 2u_{\xi\eta} \xi_x \eta_x + u_{\eta\eta} \eta_x^2 + u_{\xi} \xi_{xx} + u_{\eta} \eta_{xx}, \\ u_{xy} = u_{\xi\xi} \xi_x^2 + u_{\xi\eta} (\xi_x \eta_y + \xi_y \eta_x) + u_{\eta\eta} \eta_x \eta_y + u_{\xi} \xi_{xy} + u_{\eta} \eta_{xy}, \\ u_{yy} = u_{\xi\xi} \xi_y^2 + 2u_{\xi\eta} \xi_y \eta_y + u_{\eta\eta} \eta_y^2 + u_{\xi} \xi_{yy} + u_{\eta} \eta_{yy}. \end{array}
\]

Підставляючи ці відношення до початкового рівняння (23.3), після
простих, але громіздких обчислень отримуємо

(23.5)

\[
\overline{A}u_{\xi\xi} + \overline{B}u_{\xi\eta} + \overline{C}u_{\eta\eta} + \overline{D}u_{\xi} + \overline{E}u_{\eta} + \overline{F}u = \overline{G},
\overline{A} = A\xi_{x}^{2} + B\xi_{x}\xi_{y} + C\xi_{y}^{2},
\overline{B} = 2A\xi_{x}\eta_{y} + B(\xi_{x}\eta_{y} + \xi_{y}\eta_{x}) + 2C\xi_{y}\eta_{y},
\overline{C} = A\eta_{x}^{2} + B\eta_{x}\eta_{y} + C\eta_{y}^{2},
\overline{D} = A\xi_{xx} + B\xi_{xy} + C\xi_{yy} + D\xi_{x} + E\xi_{y},
\overline{E} = A\eta_{xx} + B\eta_{xy} + C\eta_{yy} + D\eta_{x} + E\eta_{y},
\overline{F} = F,
\overline{G} = G.
\]

Читателю предоставится возможность проделать эти выкладки при решении
задач.

\begin{verbatim}
Наш следующий шаг—выбрать функции   \xi = \xi(x, y)   и   \eta = \eta(x, y)   таким образом, чтобы коэффициенты   \overline{A}   и   \overline{C}   обратились в нуль. Это позволит нам привести исходное уравиение к канонической форме. Итак, необходимо потребовать, чтобы
\end{verbatim}

\[
\bar{A} = A\xi_x^2 + B\xi_x\xi_y + C\xi_y^2 = 0,
\]

\begin{verbatim}
\bar{C} = A\eta_x^2 + B\eta_x\eta_y + C\eta_y^2 = 0.  Эти уравнения можно представить в виде
\end{verbatim}

\[
\begin{split} &A\left[\xi_{x}/\xi_{y}\right]^{2}+B\left[\xi_{x}/\xi_{y}\right]+C=0,\\ &A\left[\eta_{x}/\eta_{y}\right]^{2}+B\left[\eta_{x}/\eta_{y}\right]^{2}+C=0. \end{split}
\]

\begin{verbatim}
Разрешив их относительно  [\xi_x/\xi_y]  и  [\eta_x/\eta_y] , получаем
\end{verbatim}

(23.7)

\[
\begin{bmatrix} \xi_x/\xi_y \end{bmatrix}
= \frac{-B + \sqrt{B^2 - 4AC}}{2A}
\]

(характеристические

\[
[\eta_x/\eta_y] = \frac{-B - \sqrt{B^2 - 4AC}}{2A}
\]

уравнения).

Кожне квадратичне рівняння для \([\xi_x/\xi_y]\) і \([\eta_x/\eta_y]\)
має два ядра, але залишаємо лише одне, щоб вони відрізнялися.

\pandocbounded{\includegraphics[keepaspectratio]{images/_page_170_Figure_3.jpeg}}

Фіг. 23.1. Характеристики \(\xi(x, y) = c\) та \(\eta(x, y) = c\) .

Задача була зведена до знаходження двох функцій \(\xi(x, y)\) і
\(\eta(x, y)\) , так що відношення \([\xi_x/\xi_y]\) і
\([\eta_x/\eta_y]\) задовольняють рівняння (23.7). Дуже легко знайти
такі функції 1) якщо звернути увагу на рис. 23.1. Щоб зрозуміти, як це
можна зробити, розглянемо просте рівняння

\[
u_{xx} - 4u_{yy} + u_x = 0.
\]

Его характеристики определяются из уравнений

\begin{verbatim}
\begin{aligned}
\frac{dy}{dx} &= -\frac{\xi_x}{\xi_y} = \frac{B - \sqrt{B^2 - 4AC}}{2A} = -2, \\
\frac{dy}{dx} &= -\frac{\eta_x}{\eta_y} = \frac{B + \sqrt{B^2 - 4AC}}{2A} = 2.
\end{aligned}
\end{verbatim}

Решая эти уравнения относительно у, получаем

\[
y = -2x + c_1,
y = 2x + c_2.
\]

\begin{verbatim}
Чтобы найти   \xi   и   \eta  , разрешим полученные уравнения относительно констант   c_1   и   c_2  , оставляя их в правой части и перенося остальные слагаемые в левую.
\end{verbatim}

Выражения, стоящие в левых частях, примем за искомые функции § и пр

\[
\xi = y + 2x = c_1,
\eta = y - 2x = c_2.
\]

1) См., например, Тихонов А. Н., Самарский А. А. Уравнения
математической физики. --- М.: Наука, 1972 и др. издания. --- Прим. ред.

Очевидно, що функції \(\xi\) і \(\eta\) введені таким чином,
задовольняють рівняння характеристик. Ці нові координати показані на
рис. 23.2. На цьому наше обговорення того, як знайти нові координати,
завершується. Останнє, що залишилося --- визначити новий тип початкового
рівняння.

\pandocbounded{\includegraphics[keepaspectratio]{images/_page_171_Figure_3.jpeg}}

Фіг. 23.2. Характеристичні координати рівняння
\(u_{xx} - 4u_{uu} + u_x = 0\) .

Це досить просто: знайти канонічну форму диференціального рівняння з
частковими похідними, підставити нові координати \(\xi(x, y)\) і
\(\eta(x,) y)\) у рівняння

\[
\overline{A}u_{\xi\xi} + \overline{B}u_{\xi\eta} + \overline{C}u_{\eta\eta} + \overline{D}u_{\xi} + \overline{E}u_{\eta} + \overline{F}u = \overline{G}
\]

,

де \(\overline{A}\) , \(\overline{B}\) , \(\overline{C}\) ,
\(\overline{D}\) , \(\overline{E}\) , \(\overline{F}\) та
\(\overline{G}\) визначаються формулами (23.6).

Прежде чем закончить лекцию, давайте на конкретном примере посмотрим,
как «работает» общий метод приведения к каноническому виду.

Приведение гиперболического уравнения \(y^2u_{xx}-x^2u_{yy}=0\) к
канонической форме

Рассмотрим уравнение

\[
y^2u_{xx}-x^2u_{yy}=0
\]

, \(x>\) , \(0 y>0\) ,

що є гіперболічним у першому квадранті. Знайдемо нові координати так,
щоб початкове рівняння набуло канонічної форми.

ШАГ 1. Решаем уравнения характеристик

\[
\frac{dy}{dx} = \frac{B + \sqrt{B^2 - 4AC}}{2A} = \frac{x}{y}
\]

(напомним, что этот шаг

\[
\frac{dy}{dx} = \frac{B - \sqrt{B^2 - 4AC}}{2A} = -\frac{x}{y}.
\]

обеспечивает \(\overline{A} = \overline{C} = 0\) ),

\pandocbounded{\includegraphics[keepaspectratio]{images/_page_172_Figure_2.jpeg}}

Фіг. 23.3. Нові координати карактора.

Інтегруючи ці звичайні диференціальні рівняння методом розділення
змінних, ми отримуємо два неявних співвідношення (за бажанням можна
виразити y через x)

\[
y^2 - x^2 = \text{const},
y^2 + x^2 = \text{const}.
\]

Следовательно, новые координаты следует вводить по формулам

\[
\xi = y^2 - x^2
\]

,

\[
\eta = y^2 + x^2  .
\]

Эти новые координаты изображены на рис. 23.3. Теперь следует получить
новое уравнение. Вычислим все коэффициенты

\begin{verbatim}
\widetilde{A} = 0   (так и должно быть; мы специально выбирали   \xi \xi   и   \eta   так, чтобы этот коэффициент обращался в нуль).
\end{verbatim}

\[
\vec{B} = 2A\xi_x\eta_x + B(\xi_x\eta_y + \xi_y\eta_x) + 2C\xi_y\eta_y = -16x^2y^2
\]

,

\begin{verbatim}
\overline{C} = 0   (по той же причине, что и   \overline{A} = 0  ),
\end{verbatim}

\[
\bar{D} = A\xi_{xx} + B\xi_{xy} + C\xi_{yy} + D\xi_x + E\xi_y = -2(x^2 + y^2)
\widetilde{E} = A\eta_{xx} + B\eta_{xy} + C\eta_{yy} + D\eta_x + E\eta_y = 2(y^2 - x^2),
\tilde{F} = F = 0
\]

,

\[
\overline{G} = G = 0
\]

,

подставим их в уравнение

\[
\overline{A}u_{\xi\xi} + \overline{B}u_{\xi\eta} + \overline{C}u_{\eta\eta} + \overline{D}u + \overline{E}u_{\eta} + \overline{F}u = \overline{G}
\]

и получим

\[
u_{\xi\eta} = \frac{-(x^2 + y^2) u_{\xi} + (y^2 - x^2) u_{\eta}}{8x^2y^3}.
\]

\begin{verbatim}
ШАГ 2. Выражая, наконец, x и y через   \xi   и   \eta  , окончательно получаем
\end{verbatim}

(23.8)

\[
u_{\xi\eta} = \frac{\eta u_{\xi} - \xi u_{\eta}}{2(\xi^2 - \eta^2)}.
\]

Задача решена. Теперь мы знаем, что нужно сделать, чтобы найти новые
координаты и привести уравнение к канонической форме.

\section{ЗАУВАЖЕННЯ}\label{ux437ux430ux443ux432ux430ux436ux435ux43dux43dux44f}

\begin{enumerate}
\def\labelenumi{\arabic{enumi}.}
\tightlist
\item
  На самом деле для гиперболических уравнений существует две
  канонические формы. Вторую можно получить из первов заменой переменных
  вида
\end{enumerate}

\[
\alpha = \alpha (\xi, \eta) = \xi + \eta,
\]

\begin{verbatim}
\beta = \beta (\xi, \eta) = \xi - \eta.  Воспользуемся этой заменой для преобразования уравнения (23.8). Так как
\end{verbatim}

\[
\begin{array}{l} u_{\xi} = u_{\alpha}\alpha_{\xi} + u_{\beta}\beta_{\xi} = u_{\alpha} + u_{\beta}. \\ u_{\eta} = u_{\alpha}\alpha_{\eta} + u_{\beta}\beta_{\eta} = u_{\alpha} - u_{\beta}, \\ u_{\xi\eta} = u_{\alpha\alpha}\alpha_{\eta} + u_{\alpha\beta}\beta_{\eta} + u_{\beta\alpha}\alpha_{\eta} + u_{\beta\beta}\beta_{\eta} = u_{\alpha\alpha} - u_{\beta\beta}, \end{array}
\]

TO

(23 9)

\[
u_{\alpha\alpha} - u_{\beta\beta} = \frac{-\beta u_{\alpha} - \alpha u_{\beta}}{2\alpha\beta}.
\]

\begin{verbatim}
При желании можно выразить   \alpha   и   \beta   через исходные координаты x и y. В результате получаем
\end{verbatim}

\[
\alpha = \xi + \eta = (y^2 - x^2) + (y^2 + x^2) = 2y^3,
\beta = \xi - \eta = (y^2 - x^2) - (y^2 + x^2) = -2x^2.
\]

\begin{itemize}
\tightlist
\item
  Читач може запитати, чому необхідно класифікувати та зводити їх до
  канонічної форми диференціальних рівнянь у похідних похідних.

  \begin{itemize}
  \tightlist
  \item
    \begin{enumerate}
    \def\labelenumi{\alph{enumi})}
    \tightlist
    \item
      Поділ рівнянь на класи гіперболічних, параболічних та еліптичних
      відповідає поділу фізичних процесів на три основні класи:
      хвильовий, дифузійний і стаціонарний. Математичні особливості
      розв'язків цих трьох типів рівнянь є абсолютно різними.
    \end{enumerate}
  \item
    \begin{enumerate}
    \def\labelenumi{\alph{enumi})}
    \setcounter{enumi}{1}
    \tightlist
    \item
      У більшості праць, присвячених розв'язанню гіперболічних задач,
      припускається, що початкове рівняння записується у канонічній
      формі, тобто у вигляді
    \end{enumerate}
  \end{itemize}
\end{itemize}

\[
u_{\xi\xi} - u_{\eta\eta} = \Psi(\xi, \, \eta, \, u_{\xi}, \, u, u_{\eta}).
\]

Якщо у нас є рівняння і ми хочемо отримати його розв'язок, ми повинні
перетворити його у канонічну форму і використати відомі результати.

\begin{enumerate}
\def\labelenumi{\alph{enumi})}
\setcounter{enumi}{2}
\tightlist
\item
  Для чисельного розв'язання канонічного гіперболічного рівняння було
  створено багато різних комп'ютерних програм.
\end{enumerate}

Функція \(\Phi\) ( \(\xi\) , \(\xi \eta\) , u, \(u_{\xi}\) ,
\(u_{\eta}\) ) вводиться у VM \(\Theta\) як підпрограма, тому рівняння
спочатку має бути канонічним. Як тільки розв'язок знайдено в нових
координатах, його завжди можна перетворити назад на початкові
координати. \# ЗАВДАННЯ - 1. Дізнайтеся, чи є наступні рівняння
гіперболічними, параболічними чи еліптичними: - : \textbf{a)}
\(u_{xx} u_{xy} = 0\) , - 6) \(u_{tt} = u_{xx} + u_x + hu\) , -
\textbf{B)} \(u_{xx} + 3u_{yy} = \sin x\) , - r)
\(u_{xx} + u_{yy} = f(x, y)\) , -
\(\mathbf{D} ) \ u_{rr} + \frac{1}{r} u_r + \frac{1}{r^2} u_{\theta \theta} = f(r, \ \theta).\)
- 2. Отримайте співвідношення (23,4), (23,5) та (23,6). - 3.
Переконайтеся, що рівняння

\[
3u_{xx} + 7u_{xy} + 2u_{yy} = 0
\]

гиперболического типа при всех x и y и найдите
\(x a p a \kappa m e p u c\) - m u + e c \kappa u + e c \kappa u + e c
\kappa u + e c \kappa u + e c \kappa u + e c \kappa u + e c \kappa u + e
c \kappa u + e c \kappa u + e c \kappa u + e c \kappa u + e c \kappa u +
e c \kappa u + e c \kappa u + e c \kappa u + e c \kappa u + e c \kappa u
+ e c \kappa u + e c \kappa u + e e c \kappa u + e c \kappa u + e c
\kappa u + e c \kappa u + e c \kappa u + e c \kappa u + e c \kappa u + e
c \kappa u + e c \kappa u + e c \kappa u + e c \kappa u + e c \kappa u +
e c \kappa u + e c \kappa u + e c \kappa u + e c \kappa u + e c \kappa u
+ e c \kappa u + e c \kappa u + e c \kap u + e c \kap

pa u + e c \kappa u + e c \kappa u + e c \kappa u + e c \kappa u + e c
\kappa u + e c \kappa u + e c \kappa u + e c \kappa u + e c \kappa u + e
c \kappa u + e c \kappa u + e c \kappa u + e c \kappa u + e c \kappa u +
e c \kappa u + e c \kappa u + e c \kappa u + e c \kappa u + e c \kappa u
+ e c \kappa u + e c \kappa u + e c \kappa u + e e c \kappa u + e c
\kappa u + e c \kappa u + e c \kappa u + e c \kappa u + e c \kappa u + e
c \kappa u + e c \kappa u + e c \kappa u + e c \kappa u + e c \kappa u +
e c \kappa u + e c \kappa u + e c \kappa u + e c \kappa u + e c \kappa u
+ e c \kappa u + e c \kappa u + e c \kappa u + e c \kappa u + e c
\kappa u + e c \kappa u + e c \kappa u + e c \kappa u + e c \kappa u + e
c \kappa u + e c \kappa u + e c \kappa u + e c \kappa u + e c \kappa u +
e c \kappa u + e c \kappa u + e c \kappa u + e c \kappa u + e c \kappa u
+ e c \kappa u + e c \kappa u + e c \kappa u + e c \kappa u + e c
\kappa u + e u + e c \kappa u + e c \kappa u + e c \kappa u + e c
\kappa u + e c \kappa u + e c \kappa u + e c \kappa u + e c \kappa u + e
c \kappa u + e c \kappa u + e c \kappa u + e c \kappa u + e c \kappa u +
e c \kappa u + e c \kappa u + e c \kappa u + e c \kappa u + e c \kappa u
+ e c \kappa u + e c ~kappa u + e c \kappa u + e c \kappa u + e c
\kappa u + e c \kappa u + e c \kappa u + e c \kappa u + e c \kappa u + e
c \kappa u + e c \kappa u + e c \kappa u + e c \kappa u + e c \kappa u +
e c \kappa u + e c \kappa u + e c \kappa u + e c \kappa u + e c \kappa u
+ e c \kappa u + e e c \kappa u + e c \kappa u + e c \kappa u + e c
\kappa u + e c \kappa u + e c \kappa u + e c \kappa u + e c \kappa u + e
c \kappa u + e c \kappa u + e c \kappa u + e c \kappa u + e c \kappa u +
e c \kappa u + e c \kappa u + e c \kappa u + e c \kappa u + e c \kappa u
+ e c \kappa u + e c \kappa u + e c \kappa u + e c \kappa u + e c
\kappa u + e c \kappa u + e c \kappa u + e c \kappa u + e c \kappa u + e
c \kappa u + e c \kappa u + e c \kappa u + e c \kappa u + e c \kappa u +
e c \kappa u + e c \kappa u + e c \kappa u + e c \kappa u + e c \kappa u
+ e c \kappa u + e c \kappa u + e c \kappa u + e u + e c \kappa u + e c
\kappa u + e c \kappa u + e c \kappa u + e c \kappa u + e c \kappa u + e
c \kappa u + e c \kappa u + e c \kappa u + e c \kappa u + e c \kappa u +
e c \kappa u + e c \kappa u + e c \kappa u + e c \kappa u + e c \kappa u
+ e c \kappa u + e c \kappa u + e c \kappa u + e c ~kappa u + e c
\kappa u + e c \kappa u + e c \kappa u + e c \kappa u + e c \kappa u + e
c \kappa u + e c \kappa u + e c \kappa u + e c \kappa u + e c \kappa u +
e c \kappa u + e c \kappa u + e c \kappa u + e c 4. Приведите уравнение
из предыдущей задачи к каноническому виду

\[
u_{\xi\eta} = \Phi(\xi, \eta, u, u_1, u_2).
\]

\begin{enumerate}
\def\labelenumi{\arabic{enumi}.}
\setcounter{enumi}{4}
\tightlist
\item
  Продолжите вадачу 4 и приведите уравнение ко второй канонической форме
\end{enumerate}

\[
u_{\alpha\alpha} - u_{\beta\beta} = \Psi(\alpha, \beta, u, u_{\alpha}, u_{\beta}).
\]

\begin{enumerate}
\def\labelenumi{\arabic{enumi}.}
\setcounter{enumi}{5}
\tightlist
\item
  Найдите характеристики уравнения
\end{enumerate}

\[
u_{xx} + 4u_{xy} = 0.
\]

Перетворіть рівняння на нові координати, розв'яжіть його і поверніться
до координат x і y, щоб отримати розв'язок початкової здогадки.

\section{ХВИЛЬОВЕ РІВНЯННЯ У ВІЛЬНОМУ ПРОСТОРІ (ДВОВИМІРНІ ТА ТРИВИМІРНІ
ЗАДАЧІ)}\label{ux445ux432ux438ux43bux44cux43eux432ux435-ux440ux456ux432ux43dux44fux43dux43dux44f-ux443-ux432ux456ux43bux44cux43dux43eux43cux443-ux43fux440ux43eux441ux442ux43eux440ux456-ux434ux432ux43eux432ux438ux43cux456ux440ux43dux456-ux442ux430-ux442ux440ux438ux432ux438ux43cux456ux440ux43dux456-ux437ux430ux434ux430ux447ux456}

МЕТАЛЕКЦІЇ: Розв'яжіть задачу з початковими умовами

(УЧП)

\[
u_{tt} = e^{2} (u_{xx} + u_{yy} + u_{zz}),
\begin{cases} -\infty < x < \infty, \\ -\infty < y < \infty, \\ -\infty < z < \infty, \end{cases}
\]

(НУ)

\[
\begin{cases} u(x, y, z, 0) = \varphi(x, y, z), \\ u_{t}(x, y, z, 0) = \psi(x, y, z) \end{cases}
\]

в трехмерном пространстве и показать, что полученное решение
удовлетворяет \emph{принципу Гюйгенса}. Для решения соответствующей
двумерной задачи

(УЧП)

\[
u_{tt} = c^2 (u_{xx} + u_{yy}), \quad
\begin{cases} -\infty < x < \infty, \\ -\infty < y < \infty, \end{cases}
\]

(НУ)

\[
Mr_1 + K = 2643
\]

використовується метод спуску, і доведено, що двовимірні розв'язки не
задовольняють принцип Гюйгенса. І нарешті, використовуючи метод спуску,
показано, що розв'язок одномірного рівняння виражається формулою
д'Аламбера, з якою ми познайомилися раніше.

У попередніх лекціях ми розглядали задачу з початковими умовами для
нескінченної вібруючої струни і показали, що її розв'язок задається
формулою д'Аламбера. Читачеві має бути зрозуміло, що в тривимірному
просторі одномірне хвильове рівняння описує плоскі хвилі. Наприклад,
звукові або електромагнітні хвилі на достатньо великих відстанях від
джерел можна вважати плоскими хвилями і, отже, описуватися одномірним
рівнянням. Існує така термінологія:

\begin{itemize}
\tightlist
\item
  \begin{enumerate}
  \def\labelenumi{\arabic{enumi})}
  \tightlist
  \item
    одномірні хвилі називаються плоскими хвилями,
  \end{enumerate}
\item
  \begin{enumerate}
  \def\labelenumi{\arabic{enumi})}
  \setcounter{enumi}{1}
  \tightlist
  \item
    Двовимірні хвилі називають циліндричними хвилями.
  \end{enumerate}
\end{itemize}

\begin{enumerate}
\def\labelenumi{\arabic{enumi})}
\setcounter{enumi}{2}
\tightlist
\item
  тривимірні хвилі називаються сферичними хвилями.
\end{enumerate}

Іншими словами, одномірне хвильове рівняння може описувати або плоскі
хвилі в просторах більшої кількості розмірів, або коливання струн.
Завдання лекції --- узагальнити формулу д'Аламбера для випадку двох і
трьох вимірів.

\section{Хвилі у тривимірному
просторі}\label{ux445ux432ux438ux43bux456-ux443-ux442ux440ux438ux432ux438ux43cux456ux440ux43dux43eux43cux443-ux43fux440ux43eux441ux442ux43eux440ux456}

Розглянемо сферичні хвилі у тривимірному просторі за заданих початкових
умов, тобто розв'яжемо задачу з початковими умовами

(24.1)

\[
\begin{cases} -\infty < x < \infty, \\ -\infty < y < \infty, \\ -\infty < z < \infty, \end{cases}
\]

(24.1)

\[
\begin{cases} u(x, y, z, 0) = \varphi(x, y, z), \\ u_{t}(x, y, z, 0) = \psi(x, y, z). \end{cases}
P_{2}
\begin{aligned} u_{tt} &= c^2 \Delta u, \\ (HY) & \begin{cases} u(x, y, z, 0) = 0, \\ u_t(x, y, z, 0) = \psi(x, y, z), \end{cases} \end{aligned}
\]

где экаком \(\Delta\) обозначен дифференциальный оператор

\begin{verbatim}
Цю задачу можна розв'язати методом  \Phi  Юр'єра  ^{1}  ), і виявляється, що розв'язок можна представити у вигляді
\end{verbatim}

\[
P_{1}
\]

\begin{verbatim}
где  \overline{\psi}  — *среднее* значение начального распределения  \psi  *по сфере* радиуса ct с центром в точке (x, y, z), т. е.
\end{verbatim}

\(+ ct \cos \varphi\) ) \((ct)^2 \sin \varphi d\varphi d\theta\) .

Аргументи функції \(\psi\) проходити по поверхні сфери, якщо значення
\(\theta\) і \(\phi\) змінюються відповідно в межах \([0, 2\pi]\) і
\([0, \pi]\) (рисунок 24.1). Цю формулу можна інтерпретувати так:

См Владимиров В. С. Уравнения математической физики. --- М.: Наука, 1971
и др. издания. --- Прим ред.

дующим образом: каждая точка пространства излучает расходящуюся (со
скоростью c) сферическую волну, через t секунд точка с координатами (x,
y, z) окажется под воздействием начального Р И С

\pandocbounded{\includegraphics[keepaspectratio]{images/_page_177_Picture_3.jpeg}}

. 24 1. Решение является средним значением начатьного распределения на
сфере.

возмущения, сосредоточенного на сфере раднуса ct с центром в заданной
точке (см. рис. 24.2).

Формула (24.3) дозволяє обчислити розв'язок задачі для більшості
початкових умов на комп'ютері. Можливо, читачу буде цікаво знайти
розв'язок для деяких простих функцій f.

\pandocbounded{\includegraphics[keepaspectratio]{images/_page_177_Picture_7.jpeg}}

Фіг. 24.2. Початкове збурення \(\psi\) поширюється з кожної точки у всіх
напрямках.

Щоб отримати повне рішення, розглянемо другу половину проблеми

(24.4)

\[
\begin{aligned} (\text{YAII}) & u_{tt} = c^2 \Delta u, \\ (\text{HY}) & \begin{cases} u = 0, \\ u_t = \varphi, \end{cases} \end{aligned}
u_h(r, \theta) = \sum_{n=0}^{\infty} r^n [a_n \cos(n\theta) + b_n \sin(n\theta)].
\Delta u = -\frac{r^2}{2} - \frac{r^3}{2} \cos{(2\theta)}
u = \frac{\partial}{\partial t} [t\overline{\varphi}].
u_p(r, \theta) = Ar^4 + Br^4 \cos(2\theta).
u(x, t) = \frac{1}{2c} \int_{x-ct}^{x+ct} \varphi(s) ds.
u_p(r, \theta) = -\frac{r^4}{32} - \frac{r^4}{24} \cos(2\theta).
u_t(x, t) = \frac{1}{2} [\varphi(x + ct) + \varphi(x - ct)].
\sum_{n=0}^{\infty} [a_n \cos(n\theta) + b_n \sin(n\theta)] - \frac{1}{32} - \frac{1}{24} \cos(2\theta) = 0.
(\text{УЧП}) \quad u_{tt} = c^2 \Delta u, \quad (x, y, z) \in \mathbb{R}^3,
u_1(r, \theta) = \frac{1}{32} + \frac{1}{24}r^2\cos(2\theta) - \frac{r^4}{24}\cos(2\theta) - \frac{r^4}{32} = \frac{(r^4 - 1)}{32} - \frac{(r^4 - r^2)}{24}\cos(2\theta).
(\text{HV}) \quad
u = u_0 + u_1 = r \cos \theta - \frac{(r^4 - 1)}{32} - \frac{(r^4 - r^2)}{24} \cos (2\theta).
\begin{cases} u = \varphi, \\ u_t = \psi \end{cases}
P_{2}
\Delta u = 0, \quad 0 < r < 1 + \frac{1}{4} \sin \theta,
\]

записывается в виде

\[
u \left( 1 + \frac{1}{4} \sin \theta, \ \theta \right) = \cos \theta, \quad 0 \le \theta \le 2\pi.
\]

\begin{verbatim}
где  \overline{\phi}  и  \overline{\psi}  — средние значения функций  \phi  и  \psi  на сфере раднусом ct с центром в точке (x, y, z).
\end{verbatim}

\begin{verbatim}
Цей розв'язок зазвичай називають формулою Пуассона для хвильового рівняння у тривимірному просторі. Це природне узагальнення формули д'Аламбера на тривимірний випадок. Найважливішим у формулі Пуассона є те, що обидва інтеграли в  \bar{\phi}  і  \bar{\psi}  взяті з поверхні сфери. Цей факт дозволяє нам надати таку важливу інтерпретацію рішення.
\end{verbatim}

\begin{verbatim}
У момент  t=t_1  розв'язок u у (x,y,z) залежить лише від початкових розподілів сфери з радіусом  ct_1  центрованим на (x,y,z) (рисунок 24.3). Тепер припустимо, що початкові розподіли  \phi  і  \psi  дорівнюють нулю всюди, крім малої сфери (див. рисунок 24.3). З часом радіус P И С
\end{verbatim}

\pandocbounded{\includegraphics[keepaspectratio]{images/_page_179_Picture_3.jpeg}}

. 24.3. Схема поясияет, как начальное возмущение влияет на точку (к, ч,
г).

\begin{verbatim}
Сфера, що оточує точку (x, y, z), зростає зі швидкістю c, так що через  t_{\hat{z}}  секунди ця сфера починає перетинати область початкового збурення, і тому розв'язок u(x, y, z, t) стає ненульовим. У  t_{\hat{z}} < t < t_{\hat{a}}  розв'язок U у (x, y, z) залишається ненульовим, оскільки сфера продовжує перетинати область початкового збурення, але на  t=t_{\hat{a}}  розв'язок знову стає нульовим. Іншими словами, розв'язок, що поширюється з початкової області збурення, має чітко визначену задню кромку. Цей загальний принцип відомий як принцип Гюйгенса для тривимірного простору. Виявляється, що передній край хвилі завжди чітко окреслений, але задній край чітко окреслений лише в просторах розмірності 3, 5, 7.... Ми вже знаємо з формули д'Аламбера, що початкові умови Р
\end{verbatim}

\pandocbounded{\includegraphics[keepaspectratio]{images/_page_185_Figure_3.jpeg}}

ис. 25.1. Графики функции \(f(x)\) = 1 и ее преобразования.

\section{Свойства
преобразований}\label{ux441ux432ux43eux439ux441ux442ux432ux430-ux43fux440ux435ux43eux431ux440ux430ux437ux43eux432ux430ux43dux438ux439}

Прежде чем приступить к решению задач, мы должны получить некоторые
свойства этих преобразований.

\begin{verbatim}
Якщо you(x, t) — функція  \partial \theta yx  змінних, і ми виконуємо перетворення змінної x, то
\end{verbatim}

(действительная форма отображения \(w=z^2\) ).

Ця форма нам знадобиться пізніше. У цьому прикладі це цікаво, оскільки
показує, як гіперболи в площині z переходять до координатних прямих u =
const і v = const у площині w.

Тепер перейдемо до вивчення одного конкретного типу функцій комплексної
змінної, які виконують конформні відображення.

\section{Определение конформного
отображения}\label{ux43eux43fux440ux435ux434ux435ux43bux435ux43dux438ux435-ux43aux43eux43dux444ux43eux440ux43cux43dux43eux433ux43e-ux43eux442ux43eux431ux440ux430ux436ux435ux43dux438ux44f}

Відображення w=\(f(z)\) комплексної площини z у комплексну площину w
називається конформним у точці \(z_{\bullet}\) площини z, якщо похідна
\(f'(z_{\bullet}) \neq 0\) . Відображення \(f(z)\) називається
конформним у регіоні D, якщо \(f'(z) \neq 0\) в кожній точці регіону D.

Наприклад, відображення \(f(z) = z^2\) конформно всюди, крім точки z =
0, оскільки \(z'(z) = 2z \neq 0\) для всіх \(z \neq 0\) . З іншого боку,
відображення \(e^z\) конформно у всій z-площині, оскільки скрізь
\(f'(z) = e^z \neq 0\) . Яке призначення конформних відображень?
Відповідь полягає в тому, що, розв'язуючи рівняння Лапласа
\(\phi_{xx} + \phi_{yy} = 0\) в певній області площини змінних x і y, ми
можемо розглядати цю площину як площину комплексної змінної z.
Розглянемо комплексне відображення w = \(f(z)\) площини z у площину w, у
якій введені координати (u, v). Відображуючи w = \(f(z\)\phi\emph{\{xy\}
+ \phi}\{uy\} = 0\()\), рівняння Лапласа
\(\varphi_{xx} + \varphi_{yy} = 0\) перетворюється на нове
диференціальне рівняння в частинних похідних, яке залежить від нових
координат і v. До --- рівняння, яке \(\phi_{xx} + \phi_{uu} = 0\)
повернеться у рівняння \(\varphi_{nn} + \varphi_{nn} = 0\) . Отже, ми
придумаємо ідею знайти конформне відображення, яке переносить область з
комплексною межею у область з простою межею (нагадаємо, що відображення
\(w=z^*\) переносить межу першого квадранта площини z у дійсну вісь
площини w).

Ось кілька прикладів, які покажуть, як можна розв'язати диференціальні
рівняння у похідних за допомогою конформних відображень.

Уравнение Лапласа в верхней полуплоскости.

Припустимо, що нам потрібно розв'язати наступну задачу Діріхле у верхній
половині видимості (рис. 47.3):

(47.1)

(Зверніть увагу, що перетворення було виконано на змінній x, і отримана
послідовність залежить лише від часу t.)

А как быть с производными? Ниже приводятся несколько полезных формул

\begin{verbatim}
\begin{aligned} (Y & \Pi) & \varphi_{xx} + \varphi_{yy} = 0, & -\infty < x < \infty, & 0 < y < \infty, \\ (47.1) & \varphi(x, 0) = \begin{cases} 0, & |x| > 1, \\ 1, & |x| \le 1. \end{cases}
\end{verbatim}

\[
w = \ln\left\{\frac{z-1}{z+1}\right\}
\begin{split} S\left[u_{t}\right] &= \frac{dS\left[u\right]}{dt}, \quad S\left[u_{tt}\right] = \frac{d^{2}S\left[u\right]}{dt^{2}}, \\ S\left[u_{xx}\right] &= -\left[n\pi/L\right]^{2}S\left[u\right] + \frac{2n\pi}{L^{2}}\left[u\left(0,\ t\right) + (-1)^{n+2}u\left(L,\ t\right)\right], \\ C\left[u_{xx}\right] &= -\left[n\pi/L\right]^{2}C\left[u\right] - \frac{2}{L}\left[u_{x}\left(0,\ t\right) + (-1)^{n+1}u_{x}\left(L,\ t\right)\right]. \end{split}
\]

\pandocbounded{\includegraphics[keepaspectratio]{images/_page_189_Picture_10.jpeg}}

Фіг. 26.1. Основна ідея принципу суперпоемації; L --- диференціальний
оператор; f --- вхід; У---виходь.

Ця базова ідея може використовуватися для розв'язання складних змішаних
задач, розбиваючи їх на прості підзавдання, які розв'язуються за
допомогою

кожне з них, а потім підсумовування отриманих результатів (звісно,
диференціальне рівняння і граничні умови мають бути лінійними).

\section{Розкладання змішаної задачі на дві
простіші}\label{ux440ux43eux437ux43aux43bux430ux434ux430ux43dux43dux44f-ux437ux43cux456ux448ux430ux43dux43eux457-ux437ux430ux434ux430ux447ux456-ux43dux430-ux434ux432ux456-ux43fux440ux43eux441ux442ux456ux448ux456}

Предположим, нам необходимо решить линейную задачу (обозначим ее P)

(УЧП)

\begin{enumerate}
\def\labelenumi{\arabic{enumi}.}
\setcounter{enumi}{17}
\tightlist
\item
\end{enumerate}

(Р) (ГУ)

\[
\begin{cases} 1-|x|, |x| < 1 \\ 0, |x| > 1 \end{cases}
\begin{cases} u(0, t) = 0, \\ u(1, t) = 0, \\ u(x, 0) = \sin(2\pi x), \quad 0 \le x \le 1. \end{cases}
\cos(ax)
\frac{2}{\pi} \frac{\sin a\omega}{\omega}
\]

Перед нами неоднородное уравнение теплопроводности, так что метод
разделения переменных неприменим. Можпо, конечно, воспользоваться
конечным синус-преобразованием по переменной x или преобразованием
Лапласа по переменной t, но вместо этого мы рассмотрим две подзадачи

\[
\frac{2}{\pi} \frac{\sin a\omega}{(\omega^2+a^2)^2}
\frac{\pi}{2} e^{-a+\omega}
\begin{cases} u(0, t) = 0, \\ u(1, t) = 0, \\ u(1, t) = 0, \end{cases}
\frac{\pi}{2} e^{-a+\omega}
\frac{\pi}{2} e^{-a+\omega}
\frac{\pi}{2} e^{-a+\omega}
\frac{\pi}{2} e^{-a+\omega}
\]

\begin{verbatim}
Каждая из этих задач легко решается и в то же время ясно, что  \mathit{suma}  решений задач  P_1  и  P_2  является решением исходной вадачи, т. е.
\end{verbatim}

\[
\frac{2ax}{(x^2+a^2)^2}
\]

Решение задачи \(P_{i}\) Розв'язання задачі \(P_i\) Розв'язання задачі
\(P_i\) Можна показати, що в загальному випадку розв'язок задачі

\[
\frac{\cos(ax), |x| > \pi/2a}{0, |x| > \pi/2a}
\frac{1-|x|, |x| < 1}{0, |x| > 1}
\begin{cases} u_t = u_{xx} + f(x, t), \\ u(0, t) = 0, \\ u(1, t) = 0, \\ u(x, 0) = \varphi(x) \end{cases}
\cos(ax)
\cos(ax)
\frac{\pi}{2} [\delta(\omega+a) + \delta(\omega-a)]
\begin{cases} u_t = u_{xx} + f(x, t), \\ u(0, t) = 0, \\ u(1, t) = 0, \\ u(x, 0) = 0 \end{cases}
f(x) = \int_{0}^{\infty} F(\omega) \sin(\omega x) d\omega
F(\omega) = \frac{?}{\pi} \int_{0}^{\infty} f(x) \sin(\omega x) dx
0 < x < \infty
\begin{cases} u_t = u_{xx}, \\ u(0, t) = 0, \\ u(1, t) = 0, \\ u(x, 0) = \varphi(x). \end{cases}
0 < \omega < x
\]

f'\,'(0)

\[
-\omega^2 F(\omega) + \frac{2}{\pi} \omega f(0)
\frac{1}{\alpha} F\left(\frac{\omega}{\alpha}\right)
\]

проследим, как оно решается методом конечного синус-преобразования. Вот
что мы должны сделать: разложить еходное воздействие \(f(x, t)\) на
компоненты, найти отклик \(U_n\) на каждую компоненту, а затем сложить
эти отклики. Математически это может быть и не совсем очевидно, но
давайте внимательно проследим за ходом решения. Разложим входящие в
уравнение с частными производными члены в ряды Фурье по синусам:

\[
\frac{2\omega}{\pi (\alpha^2 + \omega^2)}
x^{-1/2}
[2/\pi\omega]^{1/2}
\]

where

\[
\frac{2}{\pi\omega} [1-\cos(\omega a)]
\]

1

\[
e^{-1}a + \omega
\]

Зверніть увагу, по-перше, що коефіцієнти \(A_n\) , \(B_n\) \(F_n\)
дійсно є функціями часу t, оскільки ми застосували синусоїдальне
перетворення до функції двох змінних x, t, u, а по-друге, ми розклали
вхідний ефект \(f(x,t)\) на прості компоненти \(F_n\)(t) . Тепер
знайдіть відповіді \(U_n\)(t) для кожного компонента \(F_n\)(t) , потім
додайте всі відповіді \(U_n\)(t) і отримайте розв'язок \(u(x,t)\).

Щоб знайти відповіді \(U_n\)(t) , необхідно робити невеликі розрахунки з
коефіцієнтами \(A_n\)(t) і \(B_n\)(t) , щоб у вирази інтегрування була
включена лише функція u, а не її похідні \(u_i\) і \(u_{xx}\) .
Інтегруючи на частини, ми отримуємо

\[
\frac{1}{2}e^{-\alpha}\sin\omega
\frac{1-e^{-a\omega}}{\omega}
\begin{split} A_n(t) &= 2\int_0^1 u_t \sin(n\pi x) \, dx = \frac{d}{dt} \bigg[ 2\int_0^1 u(x, t) \sin(n\pi x) \, dx \bigg] = \frac{dU_n(t)}{dt} \,, \\ B_n(t) &= 2\int_0^1 u_{xx} \sin(n\pi x) \, dx = -(n\pi)^2 \, U_n(t) \,+ \\ &\quad + 2n\pi \, \big[ u(0, t) + (-1)^{n+1} \, u(1, t) \big], \end{split}
\frac{2}{\pi}F''(\omega)
f(x) = \int_{0}^{\infty} F(\omega) \cos(\omega x) d\omega \qquad \qquad F(x) = \frac{2}{\pi} \int_{0}^{\infty} f(x) \cos(\omega x) dx
0 < x < \infty \qquad \qquad 0 < \omega < \infty
\begin{cases} u(0, t) = 0, \\ u(1, t) = 0 \end{cases}
\]

j'\,'(x)

\[
S_n = \frac{2}{\pi} \int_0^{\pi} f(x) \sin(nx) dx.
\]

в последнее соотношение, получаем

\[
y = \frac{\pi}{b - a} (x - a).
\]

и, следовательно, (26.1) принимает вид

\[
f(x) = \sum_{n=1}^{\infty} S_n \sin(nx).
\]

\begin{verbatim}
Поскольку это разложение является тождеством по x, его коэффициенты должны быть равны нулю,  \tau  . е.
\end{verbatim}

\[
f(x) = \sum_{n=1}^{\infty} S_n \sin(nx)
\]

.

Таким образом, мы получили связь между входом и выходом системы, между
коэффициентами \(F_n\) и \(U_n\) . Прежде чем решать уравнения для
\(U_n\)(t) , обратимся к начальному условию

\[
S_n = \frac{2}{\pi} \int_{0}^{\pi} f(x) \sin(nx) dx
\]

Если разложить \(u(x, 0)\) в ряд Фурье по синусам и приравнять его нулю:

\[
0 \le x \le \pi
\]

\begin{verbatim}
то определяются начальные условия для функции  U_n\left(t\right)  , т. е.
\end{verbatim}

\[
n = 1, 2, ...
\]

\begin{verbatim}
,  n = 1,  2, \dots  Мы получили разложение нашей смешанной задачи на ряд простых задач об отклике системы на элементарные входные воздействия:
\end{verbatim}

(OДУ)

f'\,'(x)

(HУ) \(U_n\)(0) = 0, n = 1, 2, \ldots{} Каждое из этих уравнений легко
решается методом интегрирующего множителя (или с помощью преобразования
Лапласа). В любом случае получаем

\[
\frac{\pi}{2} e^{ax}
\]

Таким образом, мы получили отклики \(U_n\)(t) на простейшие входные
воздействия \(F_n\)(t) . Для получения решения исходной задачи осталось
совершить последний шаг---сложить все отклики

\[
C_n = \frac{2}{\pi} \int_0^{\pi} f(x) \cos(nx) dx.
\]

\begin{verbatim}
Читатель, должно быть, заметил, что каждый отклик  U_n(t)  умножается на весовой множитель  \sin(n\pi x)  . Конечно, в разложении функции f(x, t) на компоненты  F_n(t)  присутствуют те же весовые множители.
\end{verbatim}

\section{ЗАМЕЧАНИЕ}\label{ux437ux430ux43cux435ux447ux430ux43dux438ux435}

Використання скінченного синусоїдального перетворення дозволяє записати
розв'язок у вигляді ряду, але для більшості інших інтегральних
перетворень розв'язки представлені як інтеграли (неперервний розклад).
Читача можна попросити дати інтерпретацію проблеми

\[
y = \frac{\pi}{b-a}(x-a).
f(x) = \frac{C_0}{2} + \sum_{n=1}^{\infty} C_n \cos(nx).
f(x) = \frac{C_0}{2} + \sum_{n=1}^{\infty} C_n \cos(nx) \qquad C_n = \frac{2}{\pi} \int_{0}^{\pi} f(x) \cos nx \, dx
0 \le x \le \pi \qquad n = 0, 1, 2, \dots
\]

\begin{verbatim}
\begin{array}{ll} (\mathrm{Y} + \Pi) & u_{tt} = u_{xx} + \sin{(3\pi x)}, & 0 < x < 1, & 0 < t < \infty, \\ (\mathrm{\Gamma} \mathrm{Y}) & \left\{ \begin{array}{ll} u\left(0, \ t\right) = 0, \\ u\left(1, \ t\right) = 0, \\ \end{array}
\end{verbatim}

f'\,'(x)

\[
-n^{2}C_{n} - \frac{2}{\pi} \{f'(0) - (-1)^{n} f'(n)\}
\begin{array}{ll} u\left(x, \ 0\right) = \sin{(\pi x)}, \\ u_{t}\left(x, \ 0\right) = 0, \end{array}
\]

Для каждой из подзадач выберите метод решения по своему BKYCY,

\begin{enumerate}
\def\labelenumi{\arabic{enumi}.}
\setcounter{enumi}{2}
\tightlist
\item
  Нехай \(u_1\) і \(u_2\) є розв'язками наступних рівнянь:
\end{enumerate}

\begin{enumerate}
\def\labelenumi{\alph{enumi})}
\tightlist
\item
\end{enumerate}

;\\
b) \(u_t = u_{xx} + e^t\) ;

\begin{enumerate}
\def\labelenumi{\Alph{enumi})}
\setcounter{enumi}{1}
\tightlist
\item
\end{enumerate}

\begin{verbatim}
5. \cos(mx), m=1, 2, \dots
\end{verbatim}

;\\
F) \(u_t = u_{xx} + u^2\) .

\begin{verbatim}
\mathbf{\Pi}  для яких рівнянь функція  u_1 + u_2  також буде розв'язком? Які висновки ви можете зробити зі своїх відповідей?
\end{verbatim}

\begin{enumerate}
\def\labelenumi{\arabic{enumi}.}
\setcounter{enumi}{3}
\tightlist
\item
  Найдите четыре смещанные задачи, сумма решений которых дает решение
  следующей задачи:
\end{enumerate}

\[
\begin{aligned} u_t &= u_{xx} + f(x, t), & 0 < x < 1, & 0 < t < \infty, \\ & \begin{cases} u(0, t) = g_1(t), \\ u(1, t) = g_2(t), \end{cases} & 0 < t < \infty, \\ & u(x, 0) = \varphi(x), & 0 \le x \le 1. \end{aligned}
\]

\begin{enumerate}
\def\labelenumi{\arabic{enumi}.}
\setcounter{enumi}{4}
\tightlist
\item
  Решите задачу Коши
\end{enumerate}

(OAY)

x

(HY) \(U_n\)(0) = 0. Можете ли вы проверить полученное решение?

\begin{enumerate}
\def\labelenumi{\arabic{enumi}.}
\setcounter{enumi}{5}
\tightlist
\item
  Предположим, что функции \(u_1\) и \(u_2\) удовлетворяют линейным
  однородным граничным условиям
\end{enumerate}

\begin{verbatim}
\begin{cases} u_x(1, t) + h_2 u(1, t) = 0.  Чи відповідає функція u_i + u_i цим умовам? , .
\end{verbatim}

Лекция 27

\section{РІВНЯННЯ ПЕРШОГО ПОРЯДКУ (ХАРАКТЕРИСТИЧНИЙ
МЕТОД)}\label{ux440ux456ux432ux43dux44fux43dux43dux44f-ux43fux435ux440ux448ux43eux433ux43e-ux43fux43eux440ux44fux434ux43aux443-ux445ux430ux440ux430ux43aux442ux435ux440ux438ux441ux442ux438ux447ux43dux438ux439-ux43cux435ux442ux43eux434}

МЕТАЛЕКЦІЇ: Введіть поняття рівняння з похідними в похідних першого
порядку (поки що ми розглядали рівняння другого порядку) і ознайомтеся з
важливим методом розв'язання задач із початковими умовами --- методом
характеристик. Проблема, яку ми зараз розглянемо, виглядає так:

(УЧП)

\_ \textbar, -\infty \textless{} x \textless{} \infty  , 0 \textless{} t
\textless{} \infty  ,

\[
(HV)   u(x, 0) = \varphi(x), \quad -\infty < x < \infty.
\]

Слід одразу зазначити, що вперше ми розв'яжемо задачу з змінними
коефіцієнтами. Виявляється, що якщо ми рухаємося з координат (x, t) на
нові, характеристичні координати (s, \tau), то наша UFP перетвориться на
звичайне диференціальне рівняння. Тоді можна розв'язати звичайне
диференціальне рівняння, тобто знайти u(s, \tau) , і на останньому кроці
виразити s і t за x і t і отримати \(u(x, t)\).

Напомним читателю, что, когда мы решали уравнение дифф

\[
\Delta u = u_{xx} + u_{yy},   -\infty < x < \infty  ,   0 < t < \infty
\]

Константа \alpha\^{}2 виконувала роль коефіцієнта дифузії, а v ---
швидкість середовища. Очевидно, якщо \alpha=0 (тобто дифузії немає), то
розв'язок буде зсуватися як ціле число вздовж осі x зі швидкістю v
(залишається лише конвекція). Іншими словами, якщо початкова умова ---
\(u(x,0)\)=\varphi(x) , то відповідне розв'язання рівняння

\[
\Delta u = u_{rr} + \frac{1}{r}u_r + \frac{1}{r^2}u_{\theta\theta}
\]

будет иметь вид \(u(x, t)\) = \varphi(x-vt) .

Це дає нам підстави вважати, що розв'язок рівняння першого порядку

\[
\Delta u = u_{xx} + u_{yy} + u_{zz}
\]

переносится вдоль потока со скоростью

\[
\Delta u = u_{rr} + \frac{1}{r}u_r + \frac{1}{r^2}u_{\theta\theta} + u_{zz}.
\]

Звісно, якщо a і b --- константи, то розв'язок --- це хвиля, що
рухається з постійною швидкістю. Якщо \(a(x, t)\) та \(b(x, t)\)

\pandocbounded{\includegraphics[keepaspectratio]{images/_page_196_Figure_10.jpeg}}

Фіг. 27.1. Початкове значення в точці x впливає на розв'язок лише вздовж
характеристики у площині змінних (x,t). Параметр s змінюється вздовж
характеристик від нуля до нескінченності, починаючи з початкової лінії
умови; \tau --- сталий за кожною характеристикою; \tau змінюється вздовж
початкової лінії стану.

зависят от x и t, то скорость потока изменяется как вдоль потока, так и
во времени (читатель сможет увидеть, что начальная кривая сильно
искажается). Здесь так много аналогий с конвекцией!

Вернемся к нашей основной задаче:

(УЧП)

\[
\Delta u = u_{rr} + \frac{2}{r}u_r + \frac{1}{r^2}u_{\theta\theta} + \frac{\cot\theta}{r^2}
\]

(НУ) \(u(x, 0)\) = \varphi(x), -\infty \textless{} x \textless{} \infty.

Розв'язок цього рівняння першого порядку базується на наступному
фізичному факті: початкова умова в певній точці x переноситься на tx
площину вздовж прямої, що

называется характеристи-кой (см. рис. 27.1)

Таким чином, наше рівняння відрізняється від інших рівнянь (наприклад,
рівняння теплопровідності \(u_t = u_{xx})\),для яких початкове значення
в точці x впливає на розв'язок у всіх точках простору і в будь-який час.
Якщо читач пам'ятає, то початковий P I S

\pandocbounded{\includegraphics[keepaspectratio]{images/_page_197_Figure_5.jpeg}}

. 27.2. Початкове збурення \(u\left(x,0\right)\) у точці породжує дві
хвилі. Початкове збурення струни, згідно з хвильовим рівнянням,
поширюється за двома характеристиками.

Верхня струна в точці X впливає на розв'язок уздовж \часткових ліній syx
у площині TX (що відповідає двом хвилям, що рухаються) (див. рисунок
27.2).

Природно, ідея пропонує ввести дві нові координати s s і \tau (замість x
і t), щоб

\begin{itemize}
\tightlist
\item
  s змінювалися вздовж характеристичних кривих,
\item
  t змінюється вздовж початкової кривої (бажано вздовж лінії t=0).
\end{itemize}

Спочатку розглянемо нову координату s. Якщо s вибрано так, щоб
задовольнити наведену умову, то рівняння

\[
\Delta u = u_{rr} + \frac{2}{r}u_r + \frac{1}{r^2}u_{\theta\theta} + \frac{\cot\theta}{r^2}u_0 + \frac{1}{r^2\sin\theta}u_{\phi\phi}
\]

превратится в обыкновенное дифференциальное уравнение К о н е О д н

6 Воспользовавшись формулами, приведенными в тексте, решите систему
алгебраических уравнений (39.5).

Лекция 40

\section{СРАВНЕНИЕ АНАЛИТИЧЕСКИХ РЕШЕНИЙ С ЧИСЛЕННЫМИ МЕТА ЛЕКЦІЇ:
Сравнить достоинства и недостатки аналитических и численных методов
решения уравнений с частными производными. Приведена важная
математическая задача об идентификации физических величин
(параметрическая идентификация). В связи с этой задачей рассматривается
один важный пример из
биологии.}\label{ux441ux440ux430ux432ux43dux435ux43dux438ux435-ux430ux43dux430ux43bux438ux442ux438ux447ux435ux441ux43aux438ux445-ux440ux435ux448ux435ux43dux438ux439-ux441-ux447ux438ux441ux43bux435ux43dux43dux44bux43cux438-ux43cux435ux442ux430-ux43bux435ux43aux446ux456ux457-ux441ux440ux430ux432ux43dux438ux442ux44c-ux434ux43eux441ux442ux43eux438ux43dux441ux442ux432ux430-ux438-ux43dux435ux434ux43eux441ux442ux430ux442ux43aux438-ux430ux43dux430ux43bux438ux442ux438ux447ux435ux441ux43aux438ux445-ux438-ux447ux438ux441ux43bux435ux43dux43dux44bux445-ux43cux435ux442ux43eux434ux43eux432-ux440ux435ux448ux435ux43dux438ux44f-ux443ux440ux430ux432ux43dux435ux43dux438ux439-ux441-ux447ux430ux441ux442ux43dux44bux43cux438-ux43fux440ux43eux438ux437ux432ux43eux434ux43dux44bux43cux438.-ux43fux440ux438ux432ux435ux434ux435ux43dux430-ux432ux430ux436ux43dux430ux44f-ux43cux430ux442ux435ux43cux430ux442ux438ux447ux435ux441ux43aux430ux44f-ux437ux430ux434ux430ux447ux430-ux43eux431-ux438ux434ux435ux43dux442ux438ux444ux438ux43aux430ux446ux438ux438-ux444ux438ux437ux438ux447ux435ux441ux43aux438ux445-ux432ux435ux43bux438ux447ux438ux43d-ux43fux430ux440ux430ux43cux435ux442ux440ux438ux447ux435ux441ux43aux430ux44f-ux438ux434ux435ux43dux442ux438ux444ux438ux43aux430ux446ux438ux44f.-ux432-ux441ux432ux44fux437ux438-ux441-ux44dux442ux43eux439-ux437ux430ux434ux430ux447ux435ux439-ux440ux430ux441ux441ux43cux430ux442ux440ux438ux432ux430ux435ux442ux441ux44f-ux43eux434ux438ux43d-ux432ux430ux436ux43dux44bux439-ux43fux440ux438ux43cux435ux440-ux438ux437-ux431ux438ux43eux43bux43eux433ux438ux438.}

Вероятно, пришло время обсудить достоинства и недостатки аналитических и
численных решений уравнений с частными производными. Сначала выясним,
что мы подразумеваем под этими двумя типами решений.

\section{Аналитические
решения}\label{ux430ux43dux430ux43bux438ux442ux438ux447ux435ux441ux43aux438ux435-ux440ux435ux448ux435ux43dux438ux44f}

Под аналитическими решениями мы подразумеваем такие решения, в которых
неизвестная функция и выражена через независимые переменные и параметры
системы в виде формул, бесконечных рядов и интегралов.

\section{Численные
решения}\label{ux447ux438ux441ux43bux435ux43dux43dux44bux435-ux440ux435ux448ux435ux43dux438ux44f}

Під чисельним ми зрозуміємо чисельні розв'язки, отримані чисельно після
приблизної заміни початкового рівняння іншим, простішим рівнянням.
Наприклад, у методі скінченної різниці похідні виражаються через
скінченні різниці, а розв'язок диференціального рівняння в похідних у
похідних апроксимується розв'язком різницького рівняння. Результатом
такої процедури зазвичай є таблиця значень розв'язків, а для деяких
значень пояснювальних змінних.

Таблина 40.1

{\def\LTcaptype{none} % do not increment counter
\begin{longtable}[]{@{}
  >{\raggedright\arraybackslash}p{(\linewidth - 22\tabcolsep) * \real{0.1788}}
  >{\raggedright\arraybackslash}p{(\linewidth - 22\tabcolsep) * \real{0.0861}}
  >{\raggedright\arraybackslash}p{(\linewidth - 22\tabcolsep) * \real{0.1192}}
  >{\raggedright\arraybackslash}p{(\linewidth - 22\tabcolsep) * \real{0.0728}}
  >{\raggedright\arraybackslash}p{(\linewidth - 22\tabcolsep) * \real{0.0662}}
  >{\raggedright\arraybackslash}p{(\linewidth - 22\tabcolsep) * \real{0.0662}}
  >{\raggedright\arraybackslash}p{(\linewidth - 22\tabcolsep) * \real{0.0331}}
  >{\raggedright\arraybackslash}p{(\linewidth - 22\tabcolsep) * \real{0.0331}}
  >{\raggedright\arraybackslash}p{(\linewidth - 22\tabcolsep) * \real{0.0331}}
  >{\raggedright\arraybackslash}p{(\linewidth - 22\tabcolsep) * \real{0.0728}}
  >{\raggedright\arraybackslash}p{(\linewidth - 22\tabcolsep) * \real{0.1192}}
  >{\raggedright\arraybackslash}p{(\linewidth - 22\tabcolsep) * \real{0.1192}}@{}}
\toprule\noalign{}
\begin{minipage}[b]{\linewidth}\raggedright
\end{minipage} & \begin{minipage}[b]{\linewidth}\raggedright
×
\end{minipage} & \begin{minipage}[b]{\linewidth}\raggedright
\end{minipage} & \begin{minipage}[b]{\linewidth}\raggedright
\end{minipage} & \begin{minipage}[b]{\linewidth}\raggedright
\end{minipage} & \begin{minipage}[b]{\linewidth}\raggedright
\end{minipage} & \begin{minipage}[b]{\linewidth}\raggedright
\end{minipage} & \begin{minipage}[b]{\linewidth}\raggedright
\end{minipage} & \begin{minipage}[b]{\linewidth}\raggedright
\end{minipage} & \begin{minipage}[b]{\linewidth}\raggedright
\end{minipage} & \begin{minipage}[b]{\linewidth}\raggedright
\end{minipage} & \begin{minipage}[b]{\linewidth}\raggedright
\end{minipage} \\
\midrule\noalign{}
\endhead
\bottomrule\noalign{}
\endlastfoot
& 0 & 0,1 & 0,2 & 0,3 & 0,4 & 0,5 & 0,6 & 0,7 & 0,8 & 0,9 & 1 \\
00,010,020,03 & 100 & 10,20,15 & 10,34 & l: & : & : & ! & : & 10,34 &
10,20,15 & 1000 \\
; & & : & : & : & : & : & : & : & : & : & : \\
\end{longtable}
}

Теперь, когда мы знаем, что подразумевается под решением каждого из этих
типов, давайте сравним их между собой.

\section{\texorpdfstring{\textbf{Сравнительный анализ численного} и
аналитического
решений}{Сравнительный анализ численного и аналитического решений}}\label{ux441ux440ux430ux432ux43dux438ux442ux435ux43bux44cux43dux44bux439-ux430ux43dux430ux43bux438ux437-ux447ux438ux441ux43bux435ux43dux43dux43eux433ux43e-ux438-ux430ux43dux430ux43bux438ux442ux438ux447ux435ux441ux43aux43eux433ux43e-ux440ux435ux448ux435ux43dux438ux439}

Рассмотрим простейшую смещанную параболическую задачу

(40.1)

\[
y^2 - x^2 = \text{const},
\]

Розв'язок якого показано на рис. 40.1. Що краще --- мати аналітичне
рішення цієї проблеми?

(40.2)

\[
y^2 + x^2 = \text{const}.
\]

или численное, представленное при \(\alpha\) = 1 в табл. 40.1?

Вопрос хороший, а ответ на него зависит от того, что мы хотим делать с
этим решением дальше. Однако каждому из типов решений присущи очевидные
преимущества.

\section{Преимущества аналитического
решения}\label{ux43fux440ux435ux438ux43cux443ux449ux435ux441ux442ux432ux430-ux430ux43dux430ux43bux438ux442ux438ux447ux435ux441ux43aux43eux433ux43e-ux440ux435ux448ux435ux43dux438ux44f}

\begin{enumerate}
\def\labelenumi{\arabic{enumi}.}
\tightlist
\item
  Очевидно, що розв'язок (40.2) є більш інформативним, ніж таблиця
  чисел. Якщо ми хочемо обчислити розв'язок у певній точці (x, t), ми
  можемо зробити це так точно, як хочемо, просто збільшуючи кількість
  членів ряду, які розглядаються. Водночас легко отримати оцінку зверху
  щодо значення допущеної помилки.
\end{enumerate}

\[
(16.2) u_{tt} = \alpha^2 u_{xx},
u_{tt} = ku_{xx}
u_{tt} = \frac{\partial}{\partial x} [\alpha^2(x) u_x],
u_h(r, \theta) = \sum_{n=0}^{\infty} r^n [a_n \cos(n\theta) + b_n \sin(n\theta)].
u(x, 0) = f(x)
i_x + Cv_t + Gv = 0, \\ v_x + Li_t + Ri = 0,
u = \frac{\partial}{\partial t} [t\overline{\varphi}].
i_{xx} + Gv_x - CLi_{tt} - CRi_t = 0.
u_p(r, \theta) = Ar^4 + Br^4 \cos(2\theta).
v_{x} = -Li_{t} - Ri,
u(x, t) = \frac{1}{2c} \int_{x-ct}^{x+ct} \varphi(s) ds.
i_{t,s} = CLi_{tt} + (CR + GL)i_t + GRi.
u_p(r, \theta) = -\frac{r^4}{32} - \frac{r^4}{24} \cos(2\theta).
v_{xx} = CLv_{tt} + (CR + GL)v_t + GRv.
u_t(x, t) = \frac{1}{2} [\varphi(x + ct) + \varphi(x - ct)].
v_{tt} = \alpha^2 v_{xx}, \\ i_{tt} = \alpha^2 i_{xx}, \qquad \alpha^2 = 1/CL,
\sum_{n=0}^{\infty} [a_n \cos(n\theta) + b_n \sin(n\theta)] - \frac{1}{32} - \frac{1}{24} \cos(2\theta) = 0.
u_{tt} = u_{xx} - u_t
u_1(r, \theta) = \frac{1}{32} + \frac{1}{24}r^2\cos(2\theta) - \frac{r^4}{24}\cos(2\theta) - \frac{r^4}{32} = \frac{(r^4 - 1)}{32} - \frac{(r^4 - r^2)}{24}\cos(2\theta).
\]

Що ви можете сказати про точність?

\begin{enumerate}
\def\labelenumi{\arabic{enumi}.}
\setcounter{enumi}{3}
\tightlist
\item
  Як отримати послідовність випадкових точок, що лежать всередині
  трикутника T?
\end{enumerate}

РИС. 25.1. Графики функции \(f(x)\) = 1 и ее преобразования.

\section{Свойства
преобразований}\label{ux441ux432ux43eux439ux441ux442ux432ux430-ux43fux440ux435ux43eux431ux440ux430ux437ux43eux432ux430ux43dux438ux439-1}

Прежде чем приступить к решению задач, мы должны получить некоторые
свойства этих преобразований.

Якщо you(x, t) --- функція \partial \(\theta\) yx змінних, і ми
виконуємо перетворення змінної x, то

достигается на функции \$\$\(\overline{y}\)(x)\$ = \sin x\$ . Вычислите
J(\sin x) . 4. Исходя из функционала

(Зверніть увагу, що перетворення було виконано на змінній x, і отримана
послідовність залежить лише від часу t.)

А как быть с производными? Ниже приводятся несколько полезных формул

получите уравнение Эйлера --- Лагранжа

\section{ВАРИАЦИОННЫЕ МЕТОДЫ РЕШЕНИЯ УРАВНЕНИЙ С ЧАСТНЫМИ ПРОИЗВОДНЫМИ
МЕТА ЛЕКЦІЇ: Показать, как можно решать дифференциальное уравнение,
рассматривая его как уравнение Эйлера---Лагранжа для некоторого
функционала и находя некоторым новым методом функцию, минимизирующую
этот функционал. Тогда минимизирующая функция будет решением уравнения с
частными производными. Задача, конечно, состоит в том, чтобы найти такой
функционал, для которого исходное уравнение является уравнением
Эйлера---Лагранжа. Хорошо известен результат (теорема о минимуме
энергии), который гласит, что нахождение решения и некоторой
эллиптической краевой
задачи}\label{ux432ux430ux440ux438ux430ux446ux438ux43eux43dux43dux44bux435-ux43cux435ux442ux43eux434ux44b-ux440ux435ux448ux435ux43dux438ux44f-ux443ux440ux430ux432ux43dux435ux43dux438ux439-ux441-ux447ux430ux441ux442ux43dux44bux43cux438-ux43fux440ux43eux438ux437ux432ux43eux434ux43dux44bux43cux438-ux43cux435ux442ux430-ux43bux435ux43aux446ux456ux457-ux43fux43eux43aux430ux437ux430ux442ux44c-ux43aux430ux43a-ux43cux43eux436ux43dux43e-ux440ux435ux448ux430ux442ux44c-ux434ux438ux444ux444ux435ux440ux435ux43dux446ux438ux430ux43bux44cux43dux43eux435-ux443ux440ux430ux432ux43dux435ux43dux438ux435-ux440ux430ux441ux441ux43cux430ux442ux440ux438ux432ux430ux44f-ux435ux433ux43e-ux43aux430ux43a-ux443ux440ux430ux432ux43dux435ux43dux438ux435-ux44dux439ux43bux435ux440ux430ux43bux430ux433ux440ux430ux43dux436ux430-ux434ux43bux44f-ux43dux435ux43aux43eux442ux43eux440ux43eux433ux43e-ux444ux443ux43dux43aux446ux438ux43eux43dux430ux43bux430-ux438-ux43dux430ux445ux43eux434ux44f-ux43dux435ux43aux43eux442ux43eux440ux44bux43c-ux43dux43eux432ux44bux43c-ux43cux435ux442ux43eux434ux43eux43c-ux444ux443ux43dux43aux446ux438ux44e-ux43cux438ux43dux438ux43cux438ux437ux438ux440ux443ux44eux449ux443ux44e-ux44dux442ux43eux442-ux444ux443ux43dux43aux446ux438ux43eux43dux430ux43b.-ux442ux43eux433ux434ux430-ux43cux438ux43dux438ux43cux438ux437ux438ux440ux443ux44eux449ux430ux44f-ux444ux443ux43dux43aux446ux438ux44f-ux431ux443ux434ux435ux442-ux440ux435ux448ux435ux43dux438ux435ux43c-ux443ux440ux430ux432ux43dux435ux43dux438ux44f-ux441-ux447ux430ux441ux442ux43dux44bux43cux438-ux43fux440ux43eux438ux437ux432ux43eux434ux43dux44bux43cux438.-ux437ux430ux434ux430ux447ux430-ux43aux43eux43dux435ux447ux43dux43e-ux441ux43eux441ux442ux43eux438ux442-ux432-ux442ux43eux43c-ux447ux442ux43eux431ux44b-ux43dux430ux439ux442ux438-ux442ux430ux43aux43eux439-ux444ux443ux43dux43aux446ux438ux43eux43dux430ux43b-ux434ux43bux44f-ux43aux43eux442ux43eux440ux43eux433ux43e-ux438ux441ux445ux43eux434ux43dux43eux435-ux443ux440ux430ux432ux43dux435ux43dux438ux435-ux44fux432ux43bux44fux435ux442ux441ux44f-ux443ux440ux430ux432ux43dux435ux43dux438ux435ux43c-ux44dux439ux43bux435ux440ux430ux43bux430ux433ux440ux430ux43dux436ux430.-ux445ux43eux440ux43eux448ux43e-ux438ux437ux432ux435ux441ux442ux435ux43d-ux440ux435ux437ux443ux43bux44cux442ux430ux442-ux442ux435ux43eux440ux435ux43cux430-ux43e-ux43cux438ux43dux438ux43cux443ux43cux435-ux44dux43dux435ux440ux433ux438ux438-ux43aux43eux442ux43eux440ux44bux439-ux433ux43bux430ux441ux438ux442-ux447ux442ux43e-ux43dux430ux445ux43eux436ux434ux435ux43dux438ux435-ux440ux435ux448ux435ux43dux438ux44f-ux438-ux43dux435ux43aux43eux442ux43eux440ux43eux439-ux44dux43bux43bux438ux43fux442ux438ux447ux435ux441ux43aux43eux439-ux43aux440ux430ux435ux432ux43eux439-ux437ux430ux434ux430ux447ux438}

\[
w = \ln\left\{\frac{z-1}{z+1}\right\}
\]

, в области D , u = 0 , на границе D ,

эквивалентно нахождению такой функции u (тоже обращающейся в нуль на
границе области D), которая минимизирует функционал потенинальной
энергии

То есть, \(\Delta\) u = f ---уравнение Эйлера---Лагранжа для J{[}u{]}.
Функция, приближенно минимизирующая функционал J{[}u{]}, ищется методом
Ритца; тем самым мы получаем решение (приближениое) уравнения с частными
производными. В этой лекции мы познакомимся с методом Ритца и посмотрим,
как с помощью этого метода можно и и и и и и и уравнения.

Очень хорошо решать граничные задачи (как, например, задачу о
закрепленной мембране) путем отыскания гладкой поверхности, которая
минимизирует потенциальную энергию мембраны. То есть, если мы
рассматриваем исходное уравнение как уравнение Эйлера --- Лагранжа для
искоторого функционала J{[}u{]}, то решение дифференциального уравнения
можно получить путем минимизации функционала (поскольку минимизирующая
функция функционала является решением соответствующего уравнения

Эйлера --- Лагранжа). В лекции 44 мы рассматривали такие функционалы, у
которых уравнение Эйлера --- Лагранжа было обыкновенным дифференциальным
уравнением. В этой лекции рассматриваются функционалы, для которых
уравнение Эйлера --- Лагранжа относится к классу уравнений с частными
производными. Например для функционала

уравнение Эйлера---Лагранжа (подобно случаю обыкновенного
дифференциального уравнения в последней лекции) имеет вид

и, следовательно, для решения задачи Дирихле в единичном квадрате

(УЧП)

\[
(\text{УЧП}) \quad u_{tt} = u_{xx} + \sin(\pi x), \quad 0 < x < 1, \quad 0 < t < \infty,
\]

, 0 \textless{} x \textless{} 1 , 0 \textless{} y \textless{} 1 , (ГУ) u
= g , на границе квадрата,

можно, наоборот, найти такую функцию \(u(x, y)\), которая минимизирует
J{[}u{]} и равна g на границе. Вероятно, не стоит удивляться тому, что
функционал

\[
(\text{ГУ}) \quad
\]

представляет собой потенциальную эпергию мембраны и фактически мы
находим поверхность с минимальной потенциальной энерешей. Вопрос,
конечно, состоит в следующем: если дано дифференциальное уравнение, как
мы находим функционал, выражающий потенциальную энергию решения? Ответ
на этот вопрос дает хорошо известная теорема (теорема о минимуме
энергии), которая утверждает:

Решение и задачи Дирихле

\^{}\{\Lambda\}u = f в области \^{}\{D\} (ГУ) \^{}\{U\}=0 на границе
\^{}\{D\}

является той же функциен, которая минимизирует (среди функций,
удовлетворяющих граничному условию u=0) энергегический функционал

\begin{verbatim}
\begin{cases} u_{tx}, \quad 0 = \begin{cases} 1, \quad -1 < x < 1, \\ 0 \quad \text{в остальных точках,} \end{cases}
\end{verbatim}

\[
= \ln\left[\frac{z-1}{z+1}\right] + i \arg\left[\frac{z-1}{z+1}\right]
\]

и, следовательно, для решения задачи (45.1) достаточно найти функцию
\(\overline{u}\) (среди функций, обращающихся в нуль на граннце),
минимизирующую \(J\{u\}\) . Этим замечанием завершается первая часть
лекции. В оставшейся части мы покажем, как находить минимизирующую
функцию \(\overline{u}\) для J{[}u{]}, используя метод Ритца.

\section{Метод Ритца минимизации
функционалов}\label{ux43cux435ux442ux43eux434-ux440ux438ux442ux446ux430-ux43cux438ux43dux438ux43cux438ux437ux430ux446ux438ux438-ux444ux443ux43dux43aux446ux438ux43eux43dux430ux43bux43eux432}

Этот метод --- один из многих вариационных методов, описанных в
литературе по вариационному исчислению. Идея, предложенная математиком
В. Ритцем, очень проста. Метод состоит из следующих шагов.

ШАГ 1. Выбираем n и функцию u, минимизирующую функционал

\[
= \arg\left[\frac{x^2+y^2-1+i2y}{(x+1)^2+y^2}\right] =
\]

ищем в виде

\[
\begin{cases} u(x, 0) = 1, \\ u_t(x, 0) = 0, \end{cases}
\]

где функции \(\phi_1, \phi_2, \phi_3, \ldots, \phi_n\) принадлежат к
классу достаточно хороших функций и все обращаются в нуль на границе,
так что из них можно построить разумное приближение для решения задачи.
Эти функции принято называть пробными функциями. Типичный набор пробных
функций для задачи Дирихле в единичном квадрате приведен ниже:

\[
x^{2} + y^{2} = 1,
\]

--- обращаются в нуль на границе,
\(\varphi_2(x, y) = x\varphi_1(x, y),\)
\(\varphi_3(x, y) = y\varphi_1(x, y),\)
\(\varphi_4(x, y) = x^2\varphi_1(x, y),\)
\(\varphi_5(x, y) = xy\varphi_1(x, y),\)
\(\varphi_6(x, y) = y^2\varphi_1(x, y),\) \(\vdots\) \(\vdots\)

Иными словами, первые четыре приближения имеют вид

Решение задачи выполним за три следующих шага:

ШАГ І (Определение подходящего преобразования.)

Поскольку переменная x меняется от 0 до 1, мы воспользуемся конечным
преобразованием, а именно синус-преобразованием; позже станет ясно
почему. Мы могли бы решить эту задачу с помощью преобразования Лапласа
по переменной t (по трудоемкости этот метод не отличается от конечного
синус-преобразования).

ШАГ 2 (Выполнение преобразования).

Для удобства будем использовать обозначение \(S_n\)(t) = S{[}u{]} .
Применим синус-преобразование к исходному уравнению

\begin{verbatim}
\begin{aligned} & u_{ti} = c^{2}u_{xx}, & -\infty < x < +\infty, & 0 < t < \infty, \\ & (\text{НУ}) & \begin{cases} u_{t}(x, 0) = 0, \\ u_{t}(x, 0) = \begin{cases} 1, & -1 < x < 1, \\ 0, & \text{в остальных случаях.} \end{cases} \end{aligned}
\end{verbatim}

и воспользуемся тождеством для синус-преобразования, получим

в координатах (x, t).

Отже, розв'язуючи кожну з задач Дірікла всередині кола для функцій
\(u_0\) , \(u_1\) , \(u_2\) , отримуємо розв'язок

(46.9)

\begin{enumerate}
\def\labelenumi{\arabic{enumi}.}
\setcounter{enumi}{17}
\tightlist
\item
\end{enumerate}

Задачі Діріхле в деформованій області (зверніть увагу, що тут

Ми вже взяли \(\epsilon = 1/4\) :)

Як вправу, читача просять знайти пертурбацію першого порядку \(u_1\) і
перевірити, наскільки добре є приблизне розв'язання

\[
\begin{cases} 1-|x|, |x| < 1 \\ 0, |x| > 1 \end{cases}
\]

якщо розглянути рівняння з параметром

(46.11)

і побачити його розв'язок у вигляді

Підставляючи цю декомпозицію на (46.11), отримуємо таку послідовність
завдань

\[
u_{x}(0, t) = \frac{h}{T} [u(0, t) - \theta_{1}(t)]
\frac{2}{\pi} \frac{\sin a\omega}{\omega}
u_{tt} = u_{xx}, \quad 0 < x < 1, \quad 0 < t < \infty,
\frac{2}{\pi} \frac{\sin a\omega}{(\omega^2+a^2)^2}
\Delta u = 0,
q_{uu} + q_{vv} = 0
w = f(z)
\]

\begin{verbatim}
\left\{
\end{verbatim}

\[
\Rightarrow.
\frac{\pi}{2} e^{-a+\omega}
\begin{aligned} (Y & \Pi) & \varphi_{xx} + \varphi_{yy} = 0, & -\infty < x < \infty, & 0 < y < \infty, \\ (47.1) & \varphi(x, 0) = \begin{cases} 0, & |x| > 1, \\ 1, & |x| \le 1. \end{cases} \end{aligned}
\frac{\pi}{2} e^{-a+\omega}
w = \ln\left\{\frac{z-1}{z+1}\right\}
\frac{2ax}{(x^2+a^2)^2}
\begin{cases} (\text{УЧП}) & \varphi_{uu} + \varphi_{vv} = 0, \quad -\infty < u < \infty, \quad 0 < v < \pi, \\ (\text{47.2}) & \begin{cases} \varphi(u, 0) = 0, \\ \varphi(u, \pi) = 1. \end{cases} \end{cases}
\frac{\cos(ax), |x| > \pi/2a}{0, |x| > \pi/2a}
u(x, t) = \sum_{n=1}^{\infty} c_n X_n(x) T_n(t),
u_{tt} = \alpha^2 u_{xx}, \quad 0 < x < L, \quad 0 < t < \infty,
(\Gamma Y)
\]

Щоб візуально зрозуміти, як працює ця функція, ми радимо читачеві
побудувати її за кількома лініями y=c.

У другому прикладі ми перетворюємо область між двома неконцентричними
колами у кільце.

Задача Діріхле в області між двома некопоцентричними колами

Припустимо, ви хочете знайти потенціал між двома колами

\[
\cos(ax)
\quad 0 < t < \infty,
\frac{\pi}{2} [\delta(\omega+a) + \delta(\omega-a)]
\]

(HY)

\[
u = \gamma [(x-s)(x-t)-\gamma^2],
\quad 0 \le x \le L,
u(x, t) = X(x) T(t).
F(\omega) = \frac{?}{\pi} \int_{0}^{\infty} f(x) \sin(\omega x) dx
T'' - \alpha^2 \lambda T = 0,
X'' - \lambda X = 0,
u(x, t) = [C \sin(\beta x) + D \cos(\beta x)][A \sin(\alpha \beta t) + B \cos(\alpha \beta t)]
\beta_n = \frac{n\pi}{L}, \quad n = 0, 1, 2, \ldots
u_n(x, t) = X_n(x) T_n(t) = \sin(n\pi x/L) [a_n \sin(n\pi at/L) + b_n \cos(n\pi at/L)],
\]

f'\,'(0)

\[
u_n(x, t) = R_n \sin\left(\frac{n\pi x}{L}\right) \cos\left(\frac{n\pi\alpha (t - \delta_n)}{L}\right).
\frac{1}{\alpha} F\left(\frac{\omega}{\alpha}\right)
H(x-a) =
\frac{2\omega}{\pi (\alpha^2 + \omega^2)}
\]

Функція Хевісайду,

\[
[2/\pi\omega]^{1/2}
\begin{cases} 0, & x < a, \\ 1, & x \geqslant a, \end{cases}
\frac{2}{\pi\omega} [1-\cos(\omega a)]
f(x) = \mathcal{F}^{-1}[F] = \frac{1}{\sqrt{2\pi}} \int_{-\infty}^{+\infty} F(\omega) e^{-i\omega x} d\omega \qquad F(\omega) = \mathcal{F}[f] = \frac{1}{\sqrt{2\pi}} \int_{-\infty}^{+\infty} f(x) e^{-i\omega x} dx
e^{-1}a + \omega
\begin{cases} 1, |x| < a \\ 0, |x| > a \end{cases}
\frac{1}{2}e^{-\alpha}\sin\omega
\delta(x-a)
a_{n} = \frac{2}{n\pi\alpha} \int_{0}^{L} g(x) \sin(n\pi x/L) dx, \\ b_{n} = \frac{2}{L} \int_{0}^{L} f(x) \sin(n\pi x/L) dx.
u(x, t) = \sum_{n=1}^{\infty} \sin(n\pi x/L) \left[ a_n \sin(n\pi \alpha t/L) + b_n \cos(n\pi \alpha t/L) \right],
u(x, t) = \sum_{n=1}^{\infty} b_n \sin(n\pi x/L) \cos(n\pi\alpha t/L)
u(x, \theta) = f(x)
f(x) = \int_{0}^{\infty} F(\omega) \cos(\omega x) d\omega \qquad \qquad F(x) = \frac{2}{\pi} \int_{0}^{\infty} f(x) \cos(\omega x) dx
f(x) = \sum_{n=1}^{\infty} b_n \sin(n\pi x/L).
u_n(x, t) = b_n \sin(n\pi x/L) \cos(n\pi\alpha t/L).
u(x, t) = \sin(\pi x/L)\cos(\pi \alpha t/L) + 0.5\sin(3\pi x/L)\cos(3\pi \alpha t/L) + 0.25\sin(5\pi x/L)\cos(5\pi \alpha t/L).
R_n \sin(n\pi x/L) \cos[n\pi\alpha(t-\delta_n)/L],
\omega_n = \frac{n\pi\alpha}{L} = \frac{n\pi}{L} \sqrt{\frac{T}{\rho}},
S_n = \frac{2}{\pi} \int_0^{\pi} f(x) \sin(nx) dx.
u(x, 0) = \sin(\pi x/L) + 0.5 \sin(3\pi x/L),
y = \frac{\pi}{b - a} (x - a).
u(x, 0) = 0,
f(x) = \sum_{n=1}^{\infty} S_n \sin(nx).
u(x, 0) = \sin(3\pi x/L),
f(x) = \sum_{n=1}^{\infty} S_n \sin(nx)
u(x, 0) =
S_n = \frac{2}{\pi} \int_{0}^{\pi} f(x) \sin(nx) dx
\begin{cases} 2hx, & 0 \le x \le 0.5, \\ 2h(1-x), & 0.5 \le x \le 1. \end{cases}
0 \le x \le \pi
\]

К а к и

\pandocbounded{\includegraphics[keepaspectratio]{images/_page_153_Picture_8.jpeg}}

е движения будет совершать струна, если ее отпустить? 6. Решите задачу о
затухающих колебаниях струны

\[
n = 1, 2, ...
\frac{2ax}{(x^2+a^2)^2}
\]

f'\,'(x)

Представляется ли вам разумным полученное решение? Удовлетворяет ли оно
уравнению с частными производными, граничным и начальным условиям?

\begin{enumerate}
\def\labelenumi{\arabic{enumi}.}
\setcounter{enumi}{6}
\tightlist
\item
  Как вы будете решать неоднородное уравнение с заданными начальными и
  граничными условиям?
\end{enumerate}

\[
\frac{\pi}{2} e^{ax}
\cos(ax)
C_n = \frac{2}{\pi} \int_0^{\pi} f(x) \cos(nx) dx.
\frac{\pi}{2} [\delta(\omega+a) + \delta(\omega-a)]
y = \frac{\pi}{b-a}(x-a).
f(x) = \int_{0}^{\infty} F(\omega) \sin(\omega x) d\omega
f(x) = \frac{C_0}{2} + \sum_{n=1}^{\infty} C_n \cos(nx) \qquad C_n = \frac{2}{\pi} \int_{0}^{\pi} f(x) \cos nx \, dx
0 < x < \infty
\]

\begin{enumerate}
\def\labelenumi{\arabic{enumi}.}
\tightlist
\item
\end{enumerate}

\[
\frac{1}{\alpha} F\left(\frac{\omega}{\alpha}\right)
x^{-1/2}
\]

f'\,'(x)

\begin{enumerate}
\def\labelenumi{\arabic{enumi}.}
\setcounter{enumi}{4}
\tightlist
\item
  \(H(a-x)\)
\end{enumerate}

\[
\begin{array}{ll} u\left(x, \ 0\right) = \sin{(\pi x)}, \\ u_{t}\left(x, \ 0\right) = 0, \end{array}
\]

\begin{enumerate}
\def\labelenumi{\arabic{enumi}.}
\setcounter{enumi}{5}
\tightlist
\item
  \(x^{-1}\)
\end{enumerate}

\begin{enumerate}
\def\labelenumi{\alph{enumi})}
\tightlist
\item
\end{enumerate}

\begin{enumerate}
\def\labelenumi{\arabic{enumi}.}
\setcounter{enumi}{6}
\tightlist
\item
  \(\frac{x}{x^2+a^2}\)
\end{enumerate}

\begin{verbatim}
5. \cos(mx), m=1, 2, \dots
\end{verbatim}

\begin{enumerate}
\def\labelenumi{\arabic{enumi}.}
\setcounter{enumi}{8}
\tightlist
\item
  \(\arctan \frac{a}{x}\)
\end{enumerate}

\begin{verbatim}
\mathbf{\Pi}  для яких рівнянь функція  u_1 + u_2  також буде розв'язком? Які висновки ви можете зробити зі своїх відповідей?
\end{verbatim}

\begin{enumerate}
\def\labelenumi{\arabic{enumi}.}
\setcounter{enumi}{3}
\tightlist
\item
  Найдите четыре смещанные задачи, сумма решений которых дает решение
  следующей задачи:
\end{enumerate}

Откуда получаем два уравнения

\[
0 < x < \infty \qquad \qquad 0 < \omega < \infty
\]

j'\,'(x)

\begin{enumerate}
\def\labelenumi{\arabic{enumi}.}
\setcounter{enumi}{4}
\tightlist
\item
  Решите задачу Коши
\end{enumerate}

(OAY)

\[
S_n = \frac{2}{\pi} \int_0^{\pi} f(x) \sin(nx) dx.
\]

(HY) \(U_n\)(0) = 0. Можете ли вы проверить полученное решение?

\begin{enumerate}
\def\labelenumi{\arabic{enumi}.}
\setcounter{enumi}{5}
\tightlist
\item
  Предположим, что функции \(u_1\) и \(u_2\) удовлетворяют линейным
  однородным граничным условиям
\end{enumerate}

\[
y = \frac{\pi}{b - a} (x - a).
\]

\begin{verbatim}
\begin{cases}  u_x(1, t) + h_2 u(1, t) = 0.  Будет ли функция   u_i + u_i   удовлетворять этим условиям? , .
\end{verbatim}

Лекция 27

\section{УРАВНЕНИЯ ПЕРВОГО ПОРЯДКА (МЕТОД
ХАРАКТЕРИСТИК)}\label{ux443ux440ux430ux432ux43dux435ux43dux438ux44f-ux43fux435ux440ux432ux43eux433ux43e-ux43fux43eux440ux44fux434ux43aux430-ux43cux435ux442ux43eux434-ux445ux430ux440ux430ux43aux442ux435ux440ux438ux441ux442ux438ux43a}

ЦЕЛЬ ЛЕКЦИИ: Ввести понятие уравнения с частными производными первого
порядка (до сих пор мы рассматривали уравнения второго порядка) и
познакомиться с важным методом решения задач с начальными
условиями---методом характеристик. Задача, которой мы теперь займемся,
имеет вид

\begin{verbatim}
(УЧП) \end{cases}
\end{verbatim}

\[
f(x) = \sum_{n=1}^{\infty} S_n \sin(nx).
S_n = \frac{2}{\pi} \int_{0}^{\pi} f(x) \sin(nx) dx
\]

Сразу же отметим, что впервые мы будем решать задачу с переменными
коэффициентами. Оказывается, если от координат (x, t) перейти к новым,
характеристическим координатам (s, \tau) , то наше УЧП превратится в
обыкновенное дифференциальное уравнение. Затем можно решить обыкновенное
дифференциальное уравнение, т. е. найти u(s, \tau) , а на последнем шаге
выразить s и t через x и t и получить \(u(x, t)\).

Напомним читателю, что, когда мы решали уравнение дифф

Воспользовавшись ортогональностью семейства функции (sin (nnx)) на
отреже {[}0, 1{]}, находим формулы для коэффициентов

(21.6)

константа \alpha\^{}2 играла роль коэффициента диффузии, а v---была
скоростью среды. Очевидно, что если \alpha=0 (т. е. диффузии нет), то
решение будет смещаться как целое вдоль оси x со скоростью v (у нас
осталась только конвекция). Другими словами, если начальное условие
\(u(x,0)\)=\varphi(x) , то соответствующее решение уравнения

\[
\frac{\pi}{2} e^{ax}
\]

будет иметь вид \(u(x, t)\) = \varphi(x-vt) .

Это дает нам основания думать, что решение уравнения первого порядка

\[
C_n = \frac{2}{\pi} \int_0^{\pi} f(x) \cos(nx) dx.
\]

переносится вдоль потока со скоростью

\[
y = \frac{\pi}{b-a}(x-a).
\]

Конечно, если a и b---константы, то решение представляет собой бегущую с
постоянной скоростью волну. Если же \(a(x, t)\) и \(b(x, t)\)

\[
f(x) = \frac{C_0}{2} + \sum_{n=1}^{\infty} C_n \cos(nx).
\]

РИС. 27.1. Начальное значение в точке x оказывает влияние на решение
только вдоль характеристики в плоскости переменных (x,t). Параметр s
изменяется вдоль характеристик от нуля до бесконечности, начиная с линии
начального условия; \tau   постоянна вдоль каждой характеристики;
\tau   изменяется вдоль линии начального условия.

зависят от x и t, то скорость потока изменяется как вдоль потока, так и
во времени (читатель сможет увидеть, что начальная кривая сильно
искажается). Здесь так много аналогий с конвекцией!

Вернемся к нашей основной задаче:

(УЧП)

\[
0 \le x \le \pi \qquad n = 0, 1, 2, \dots
\]

(НУ) \(u(x, 0)\) = \varphi(x), -\infty \textless{} x \textless{} \infty.

Решение этого линейного уравнения первого порядка основывается на
следующем физическом факте: начальное условие в некоторой точке x
переносится в tx-плоскости вдоль линии, которая

называется характеристи-кой (см. рис. 27.1)

Этим наше уразнение отличается от других уравнений (таких, как уравнение
теплопроводности \(u_t = u_{xx} )\),для которых начальное значение в
точке x оказывает влияние на решение во всех точках пространства и во
все моменты времени. Если читатель помнит, то начальное сме- Р И С

f'\,'(x)

. 27.2. Начальное возмущение \(u\left(x,0\right)\) в точке порождает две
волны. Начальное возмущение струны согласно волновому уравнению
распространяется вдоль двух характеристик.

щение скрипичной струны в точке x оказывает влияние на решение вдоль
\partial syx линий в tx-плоскости (соответствующих двум бегущим волнам)
(см. рис. 27.2).

Естественно, напрашивается идея ввести две новые координаты s и
\tau   (вместо x и t) так, чтобы

\begin{itemize}
\tightlist
\item
  s изменялась вдоль характеристических кривых,
\item
  т изменялась вдоль начальной кривой (желательно вдоль линии t=0).
\end{itemize}

Рассмотрим сначала новую координату s. Если s выбрать так, что она
удовлетворяет приведенному выше условию, то уравнение

\begin{verbatim}
2.  \frac{a_{0}}{2} + \sum_{n=1}^{\infty} a_{n} \cos(nx)   a_{n}  3.  f(n-x)   (-1)^{n} \frac{2}{\pi} C_{n}  4. 1
\end{verbatim}

превратится в обыкновенное дифференциальное уравнение К о н

\begin{verbatim}
5.   \cos(mx), m=1, 2, \dots 
\end{verbatim}

ечно, остается открытым вопрос, как найти эти характеристики. Ответ
прост: мы выберем характеристики \{{[}\(x(s)\), \(t(s)\){]}: 0
\textless{} s \textless{} \infty\} так, чтобы они удовлетворяли системе

\$\$

\begin{verbatim}
\begin{cases}
\end{verbatim}

\end{document}
